% ============================================================
% Hugging Face Daily Papers 半月综述：2026年5月1日--5月14日
% ============================================================

\section{Hugging Face Daily Papers 半月综述：2026年5月1日--5月14日}

\begin{conceptbox}
\textbf{数据来源}：\href{https://huggingface.co/papers/month/2026-05}{Hugging Face Papers 2026 年 5 月月度页面}；\href{https://huggingface.co/papers}{Hugging Face 每日论文首页}；每个条目对应 arXiv 摘要页。\\
\textbf{整理日期}：2026-05-14\\
\textbf{覆盖范围}：2026-05-01 至 2026-05-14 月度页面中的 105 篇每日论文。\\
\textbf{整理原则}：按问题域聚类，而不是按日期机械罗列；每篇保留动机、背景、方法、洞察、结论；每类之后给出趋势和下一阶段预测。
\end{conceptbox}

我的总体判断：五月上半月的 AI 论文主线正在从“单个模型能力展示”转向“可行动、可搜索、可记忆、可奖励、可评测、可部署的系统闭环”。

\subsection{六大类总览}
\begin{enumerate}[leftmargin=1.5em]
  \item \textbf{具身智能、视觉—语言—动作模型与世界行动模型}：12 篇。本类关注模型如何从视觉/语言理解走向可执行动作、可预测世界和可交互的三维、四维状态。
  \item \textbf{智能体、搜索、记忆与科研自动化}：29 篇。本类关注智能体如何搜索、记忆、调用工具、构造技能并自动化科研工作流。
  \item \textbf{后训练、可验证奖励强化学习、奖励建模与测试时计算}：24 篇。本类关注奖励、策略优化、在线策略数据演化、评价细则和测试时计算如何改变模型能力。
  \item \textbf{统一多模态理解生成、图像与视频生成}：19 篇。本类关注图像、视频、音视觉和统一多模态模型如何从生成质量走向可控编辑与统一交互。
  \item \textbf{基础架构、表示学习与推理效率}：15 篇。本类关注潜空间、词元、注意力、预填充、模型合并和超深网络稳定性等底层问题。
  \item \textbf{评测、安全、隐私与可信科学工具}：6 篇。本类关注基准、隐私、红队、可信推理和科学工具如何为能力增长提供约束。
\end{enumerate}

\subsection{具身智能、视觉—语言—动作模型与世界行动模型}

本类关注模型如何从视觉/语言理解走向可执行动作、可预测世界和可交互的三维、四维状态。

\subsubsection{论文卡片}
\begin{enumerate}[leftmargin=1.5em]
  \item \textbf{\href{https://huggingface.co/papers/2605.12090}{世界行动模型：具身智能的下一前沿}}（\href{https://arxiv.org/abs/2605.12090}{2605.12090}；每日论文日期：2026-05-13；点赞数：51）
  \begin{itemize}[leftmargin=1.5em]
    \item \textbf{动机}：让模型不仅“看懂世界”，还要能预测行动后果，并把动作可靠地落到真实或仿真环境中。
    \item \textbf{背景}：视觉—语言—动作模型、视频生成、三维/四维表征和世界模型正在汇合，但现有系统常缺少显式动力学、反事实预测和真实部署鲁棒性。
    \item \textbf{方法}：论文通常通过新的世界模型、动作推理模型、三维/四维生成框架、物理约束或部署策略，连接感知、状态转移和动作执行。本文题目中的核心对象是“世界行动模型：具身智能的下一前沿”，可把它理解为该方向的一种具体实现或问题重述。
    \item \textbf{洞察}：真正有用的世界模型必须同时回答状态会如何变化、动作是否可行以及失败风险在哪里。
    \item \textbf{结论}：它的结论价值在于推动世界模型从“生成画面”走向“支持行动决策”。
  \end{itemize}
  \item \textbf{\href{https://huggingface.co/papers/2605.09131}{MCP-Cosmos：在 MCP 环境中用世界模型增强复杂任务执行智能体}}（\href{https://arxiv.org/abs/2605.09131}{2605.09131}；每日论文日期：2026-05-13；点赞数：27）
  \begin{itemize}[leftmargin=1.5em]
    \item \textbf{动机}：让模型不仅“看懂世界”，还要能预测行动后果，并把动作可靠地落到真实或仿真环境中。
    \item \textbf{背景}：视觉—语言—动作模型、视频生成、三维/四维表征和世界模型正在汇合，但现有系统常缺少显式动力学、反事实预测和真实部署鲁棒性。
    \item \textbf{方法}：论文通常通过新的世界模型、动作推理模型、三维/四维生成框架、物理约束或部署策略，连接感知、状态转移和动作执行。本文题目中的核心对象是“MCP-Cosmos：在 MCP 环境中用世界模型增强复杂任务执行智能体”，可把它理解为该方向的一种具体实现或问题重述。
    \item \textbf{洞察}：真正有用的世界模型必须同时回答状态会如何变化、动作是否可行以及失败风险在哪里。
    \item \textbf{结论}：它的结论价值在于推动世界模型从“生成画面”走向“支持行动决策”。
  \end{itemize}
  \item \textbf{\href{https://huggingface.co/papers/2605.10922}{Pixal3D：从图像生成像素对齐的 3D 内容}}（\href{https://arxiv.org/abs/2605.10922}{2605.10922}；每日论文日期：2026-05-12；点赞数：22）
  \begin{itemize}[leftmargin=1.5em]
    \item \textbf{动机}：让模型不仅“看懂世界”，还要能预测行动后果，并把动作可靠地落到真实或仿真环境中。
    \item \textbf{背景}：视觉—语言—动作模型、视频生成、三维/四维表征和世界模型正在汇合，但现有系统常缺少显式动力学、反事实预测和真实部署鲁棒性。
    \item \textbf{方法}：论文通常通过新的世界模型、动作推理模型、三维/四维生成框架、物理约束或部署策略，连接感知、状态转移和动作执行。本文题目中的核心对象是“Pixal3D：从图像生成像素对齐的 3D 内容”，可把它理解为该方向的一种具体实现或问题重述。
    \item \textbf{洞察}：真正有用的世界模型必须同时回答状态会如何变化、动作是否可行以及失败风险在哪里。
    \item \textbf{结论}：它的结论价值在于推动世界模型从“生成画面”走向“支持行动决策”。
  \end{itemize}
  \item \textbf{\href{https://huggingface.co/papers/2605.06222}{何时信任想象：面向世界行动模型的自适应动作执行}}（\href{https://arxiv.org/abs/2605.06222}{2605.06222}；每日论文日期：2026-05-08；点赞数：39）
  \begin{itemize}[leftmargin=1.5em]
    \item \textbf{动机}：让模型不仅“看懂世界”，还要能预测行动后果，并把动作可靠地落到真实或仿真环境中。
    \item \textbf{背景}：视觉—语言—动作模型、视频生成、三维/四维表征和世界模型正在汇合，但现有系统常缺少显式动力学、反事实预测和真实部署鲁棒性。
    \item \textbf{方法}：论文通常通过新的世界模型、动作推理模型、三维/四维生成框架、物理约束或部署策略，连接感知、状态转移和动作执行。本文题目中的核心对象是“何时信任想象：面向世界行动模型的自适应动作执行”，可把它理解为该方向的一种具体实现或问题重述。
    \item \textbf{洞察}：真正有用的世界模型必须同时回答状态会如何变化、动作是否可行以及失败风险在哪里。
    \item \textbf{结论}：它的结论价值在于推动世界模型从“生成画面”走向“支持行动决策”。
  \end{itemize}
  \item \textbf{\href{https://huggingface.co/papers/2605.03269}{RLDX-1 技术报告}}（\href{https://arxiv.org/abs/2605.03269}{2605.03269}；每日论文日期：2026-05-07；点赞数：117）
  \begin{itemize}[leftmargin=1.5em]
    \item \textbf{动机}：让模型不仅“看懂世界”，还要能预测行动后果，并把动作可靠地落到真实或仿真环境中。
    \item \textbf{背景}：视觉—语言—动作模型、视频生成、三维/四维表征和世界模型正在汇合，但现有系统常缺少显式动力学、反事实预测和真实部署鲁棒性。
    \item \textbf{方法}：论文通常通过新的世界模型、动作推理模型、三维/四维生成框架、物理约束或部署策略，连接感知、状态转移和动作执行。本文题目中的核心对象是“RLDX-1 技术报告”，可把它理解为该方向的一种具体实现或问题重述。
    \item \textbf{洞察}：这篇工作的关键不在单点技巧，而在它把数据、模型、推理、奖励与评测中的某个环节重新组织进系统闭环。
    \item \textbf{结论}：它的结论价值在于推动世界模型从“生成画面”走向“支持行动决策”。
  \end{itemize}
  \item \textbf{\href{https://huggingface.co/papers/2604.28196}{HERMES++：面向 3D 场景理解与生成的统一驾驶世界模型}}（\href{https://arxiv.org/abs/2604.28196}{2604.28196}；每日论文日期：2026-05-07；点赞数：71）
  \begin{itemize}[leftmargin=1.5em]
    \item \textbf{动机}：让模型不仅“看懂世界”，还要能预测行动后果，并把动作可靠地落到真实或仿真环境中。
    \item \textbf{背景}：视觉—语言—动作模型、视频生成、三维/四维表征和世界模型正在汇合，但现有系统常缺少显式动力学、反事实预测和真实部署鲁棒性。
    \item \textbf{方法}：论文通常通过新的世界模型、动作推理模型、三维/四维生成框架、物理约束或部署策略，连接感知、状态转移和动作执行。本文题目中的核心对象是“HERMES++：面向 3D 场景理解与生成的统一驾驶世界模型”，可把它理解为该方向的一种具体实现或问题重述。
    \item \textbf{洞察}：真正有用的世界模型必须同时回答状态会如何变化、动作是否可行以及失败风险在哪里。
    \item \textbf{结论}：它的结论价值在于推动世界模型从“生成画面”走向“支持行动决策”。
  \end{itemize}
  \item \textbf{\href{https://huggingface.co/papers/2605.05163}{PhysForge：为交互式虚拟世界生成物理扎根的 3D 资产}}（\href{https://arxiv.org/abs/2605.05163}{2605.05163}；每日论文日期：2026-05-07；点赞数：35）
  \begin{itemize}[leftmargin=1.5em]
    \item \textbf{动机}：让模型不仅“看懂世界”，还要能预测行动后果，并把动作可靠地落到真实或仿真环境中。
    \item \textbf{背景}：视觉—语言—动作模型、视频生成、三维/四维表征和世界模型正在汇合，但现有系统常缺少显式动力学、反事实预测和真实部署鲁棒性。
    \item \textbf{方法}：论文通常通过新的世界模型、动作推理模型、三维/四维生成框架、物理约束或部署策略，连接感知、状态转移和动作执行。本文题目中的核心对象是“PhysForge：为交互式虚拟世界生成物理扎根的 3D 资产”，可把它理解为该方向的一种具体实现或问题重述。
    \item \textbf{洞察}：真正有用的世界模型必须同时回答状态会如何变化、动作是否可行以及失败风险在哪里。
    \item \textbf{结论}：它的结论价值在于推动世界模型从“生成画面”走向“支持行动决策”。
  \end{itemize}
  \item \textbf{\href{https://huggingface.co/papers/2605.02881}{MolmoAct2：面向真实部署的动作推理模型}}（\href{https://arxiv.org/abs/2605.02881}{2605.02881}；每日论文日期：2026-05-05；点赞数：287）
  \begin{itemize}[leftmargin=1.5em]
    \item \textbf{动机}：让模型不仅“看懂世界”，还要能预测行动后果，并把动作可靠地落到真实或仿真环境中。
    \item \textbf{背景}：视觉—语言—动作模型、视频生成、三维/四维表征和世界模型正在汇合，但现有系统常缺少显式动力学、反事实预测和真实部署鲁棒性。
    \item \textbf{方法}：论文通常通过新的世界模型、动作推理模型、三维/四维生成框架、物理约束或部署策略，连接感知、状态转移和动作执行。本文题目中的核心对象是“MolmoAct2：面向真实部署的动作推理模型”，可把它理解为该方向的一种具体实现或问题重述。
    \item \textbf{洞察}：真正有用的世界模型必须同时回答状态会如何变化、动作是否可行以及失败风险在哪里。
    \item \textbf{结论}：它的结论价值在于推动世界模型从“生成画面”走向“支持行动决策”。
  \end{itemize}
  \item \textbf{\href{https://huggingface.co/papers/2605.00781}{Map2World：由分割图约束的文本到 3D 世界生成}}（\href{https://arxiv.org/abs/2605.00781}{2605.00781}；每日论文日期：2026-05-04；点赞数：25）
  \begin{itemize}[leftmargin=1.5em]
    \item \textbf{动机}：让模型不仅“看懂世界”，还要能预测行动后果，并把动作可靠地落到真实或仿真环境中。
    \item \textbf{背景}：视觉—语言—动作模型、视频生成、三维/四维表征和世界模型正在汇合，但现有系统常缺少显式动力学、反事实预测和真实部署鲁棒性。
    \item \textbf{方法}：论文通常通过新的世界模型、动作推理模型、三维/四维生成框架、物理约束或部署策略，连接感知、状态转移和动作执行。本文题目中的核心对象是“Map2World：由分割图约束的文本到 3D 世界生成”，可把它理解为该方向的一种具体实现或问题重述。
    \item \textbf{洞察}：真正有用的世界模型必须同时回答状态会如何变化、动作是否可行以及失败风险在哪里。
    \item \textbf{结论}：它的结论价值在于推动世界模型从“生成画面”走向“支持行动决策”。
  \end{itemize}
  \item \textbf{\href{https://huggingface.co/papers/2604.28185}{新时代的视觉生成：从原子映射到智能体式世界建模}}（\href{https://arxiv.org/abs/2604.28185}{2604.28185}；每日论文日期：2026-05-01；点赞数：89）
  \begin{itemize}[leftmargin=1.5em]
    \item \textbf{动机}：让模型不仅“看懂世界”，还要能预测行动后果，并把动作可靠地落到真实或仿真环境中。
    \item \textbf{背景}：视觉—语言—动作模型、视频生成、三维/四维表征和世界模型正在汇合，但现有系统常缺少显式动力学、反事实预测和真实部署鲁棒性。
    \item \textbf{方法}：论文通常通过新的世界模型、动作推理模型、三维/四维生成框架、物理约束或部署策略，连接感知、状态转移和动作执行。本文题目中的核心对象是“新时代的视觉生成：从原子映射到智能体式世界建模”，可把它理解为该方向的一种具体实现或问题重述。
    \item \textbf{洞察}：真正有用的世界模型必须同时回答状态会如何变化、动作是否可行以及失败风险在哪里。
    \item \textbf{结论}：它的结论价值在于推动世界模型从“生成画面”走向“支持行动决策”。
  \end{itemize}
  \item \textbf{\href{https://huggingface.co/papers/2604.27711}{ExoActor：把外视角视频生成用于可泛化交互式人形控制}}（\href{https://arxiv.org/abs/2604.27711}{2604.27711}；每日论文日期：2026-05-01；点赞数：41）
  \begin{itemize}[leftmargin=1.5em]
    \item \textbf{动机}：让模型不仅“看懂世界”，还要能预测行动后果，并把动作可靠地落到真实或仿真环境中。
    \item \textbf{背景}：视觉—语言—动作模型、视频生成、三维/四维表征和世界模型正在汇合，但现有系统常缺少显式动力学、反事实预测和真实部署鲁棒性。
    \item \textbf{方法}：论文通常通过新的世界模型、动作推理模型、三维/四维生成框架、物理约束或部署策略，连接感知、状态转移和动作执行。本文题目中的核心对象是“ExoActor：把外视角视频生成用于可泛化交互式人形控制”，可把它理解为该方向的一种具体实现或问题重述。
    \item \textbf{洞察}：生成模型的核心竞争正从单帧质量转向可控性、一致性和工作流可用性。
    \item \textbf{结论}：它的结论价值在于推动世界模型从“生成画面”走向“支持行动决策”。
  \end{itemize}
  \item \textbf{\href{https://huggingface.co/papers/2604.28130}{MoCapAnything V2：面向任意骨架的端到端动作捕捉}}（\href{https://arxiv.org/abs/2604.28130}{2604.28130}；每日论文日期：2026-05-01；点赞数：22）
  \begin{itemize}[leftmargin=1.5em]
    \item \textbf{动机}：让模型不仅“看懂世界”，还要能预测行动后果，并把动作可靠地落到真实或仿真环境中。
    \item \textbf{背景}：视觉—语言—动作模型、视频生成、三维/四维表征和世界模型正在汇合，但现有系统常缺少显式动力学、反事实预测和真实部署鲁棒性。
    \item \textbf{方法}：论文通常通过新的世界模型、动作推理模型、三维/四维生成框架、物理约束或部署策略，连接感知、状态转移和动作执行。本文题目中的核心对象是“MoCapAnything V2：面向任意骨架的端到端动作捕捉”，可把它理解为该方向的一种具体实现或问题重述。
    \item \textbf{洞察}：这篇工作的关键不在单点技巧，而在它把数据、模型、推理、奖励与评测中的某个环节重新组织进系统闭环。
    \item \textbf{结论}：它的结论价值在于推动世界模型从“生成画面”走向“支持行动决策”。
  \end{itemize}
\end{enumerate}

\subsubsection{类别趋势}
趋势：具身智能正在从单任务策略转向开放动作模型、动作词元化、世界状态预测和三维/四维空间表征的一体化。

\subsubsection{下一阶段预测}
预测：下一阶段会出现更多世界—行动基础模型，评测重点会从视觉逼真度转向动作成功率、反事实预测、风险感知和可交互仿真。

\subsection{智能体、搜索、记忆与科研自动化}

本类关注智能体如何搜索、记忆、调用工具、构造技能并自动化科研工作流。

\subsubsection{论文卡片}
\begin{enumerate}[leftmargin=1.5em]
  \item \textbf{\href{https://huggingface.co/papers/2605.09530}{MemPrivacy：面向边云智能体的隐私保护个性化记忆管理}}（\href{https://arxiv.org/abs/2605.09530}{2605.09530}；每日论文日期：2026-05-13；点赞数：128）
  \begin{itemize}[leftmargin=1.5em]
    \item \textbf{动机}：让模型从一次性问答转向长期任务执行，能够主动检索、保存记忆、组合工具并生成可复用技能。
    \item \textbf{背景}：传统检索增强生成与单轮工具调用难以支撑开放任务，长期记忆、隐私、轨迹质量和科研过程建模成为新瓶颈。
    \item \textbf{方法}：论文通常通过高质量轨迹、工具编排、长期记忆结构、多智能体协作或科研图谱，增强智能体的任务闭环。本文题目中的核心对象是“MemPrivacy：面向边云智能体的隐私保护个性化记忆管理”，可把它理解为该方向的一种具体实现或问题重述。
    \item \textbf{洞察}：长期记忆不是简单追加上下文，而是需要压缩、权限、隐私和可更新状态共同设计。
    \item \textbf{结论}：它的结论价值在于把智能体能力从提示词技巧推进到可训练、可评估的系统流程。
  \end{itemize}
  \item \textbf{\href{https://huggingface.co/papers/2605.12357}{δ-mem：面向大语言模型的高效在线记忆}}（\href{https://arxiv.org/abs/2605.12357}{2605.12357}；每日论文日期：2026-05-13；点赞数：89）
  \begin{itemize}[leftmargin=1.5em]
    \item \textbf{动机}：让模型从一次性问答转向长期任务执行，能够主动检索、保存记忆、组合工具并生成可复用技能。
    \item \textbf{背景}：传统检索增强生成与单轮工具调用难以支撑开放任务，长期记忆、隐私、轨迹质量和科研过程建模成为新瓶颈。
    \item \textbf{方法}：论文通常通过高质量轨迹、工具编排、长期记忆结构、多智能体协作或科研图谱，增强智能体的任务闭环。本文题目中的核心对象是“δ-mem：面向大语言模型的高效在线记忆”，可把它理解为该方向的一种具体实现或问题重述。
    \item \textbf{洞察}：长期记忆不是简单追加上下文，而是需要压缩、权限、隐私和可更新状态共同设计。
    \item \textbf{结论}：它的结论价值在于把智能体能力从提示词技巧推进到可训练、可评估的系统流程。
  \end{itemize}
  \item \textbf{\href{https://huggingface.co/papers/2605.10899}{RubricEM：超越可验证奖励的评价细则引导策略分解元强化学习}}（\href{https://arxiv.org/abs/2605.10899}{2605.10899}；每日论文日期：2026-05-13；点赞数：66）
  \begin{itemize}[leftmargin=1.5em]
    \item \textbf{动机}：让模型从一次性问答转向长期任务执行，能够主动检索、保存记忆、组合工具并生成可复用技能。
    \item \textbf{背景}：传统检索增强生成与单轮工具调用难以支撑开放任务，长期记忆、隐私、轨迹质量和科研过程建模成为新瓶颈。
    \item \textbf{方法}：论文通常通过高质量轨迹、工具编排、长期记忆结构、多智能体协作或科研图谱，增强智能体的任务闭环。本文题目中的核心对象是“RubricEM：超越可验证奖励的评价细则引导策略分解元强化学习”，可把它理解为该方向的一种具体实现或问题重述。
    \item \textbf{洞察}：奖励信号越能贴近任务结构，后训练越可能稳定地产生能力增益。
    \item \textbf{结论}：它的结论价值在于把智能体能力从提示词技巧推进到可训练、可评估的系统流程。
  \end{itemize}
  \item \textbf{\href{https://huggingface.co/papers/2605.12481}{ToolCUA：面向计算机使用智能体的图形界面—工具路径最优编排}}（\href{https://arxiv.org/abs/2605.12481}{2605.12481}；每日论文日期：2026-05-13；点赞数：23）
  \begin{itemize}[leftmargin=1.5em]
    \item \textbf{动机}：让模型从一次性问答转向长期任务执行，能够主动检索、保存记忆、组合工具并生成可复用技能。
    \item \textbf{背景}：传统检索增强生成与单轮工具调用难以支撑开放任务，长期记忆、隐私、轨迹质量和科研过程建模成为新瓶颈。
    \item \textbf{方法}：论文通常通过高质量轨迹、工具编排、长期记忆结构、多智能体协作或科研图谱，增强智能体的任务闭环。本文题目中的核心对象是“ToolCUA：面向计算机使用智能体的图形界面—工具路径最优编排”，可把它理解为该方向的一种具体实现或问题重述。
    \item \textbf{洞察}：这篇工作的关键不在单点技巧，而在它把数据、模型、推理、奖励与评测中的某个环节重新组织进系统闭环。
    \item \textbf{结论}：它的结论价值在于把智能体能力从提示词技巧推进到可训练、可评估的系统流程。
  \end{itemize}
  \item \textbf{\href{https://huggingface.co/papers/2605.10832}{迈向视觉原生多模态深度搜索智能体的在线策略数据演化}}（\href{https://arxiv.org/abs/2605.10832}{2605.10832}；每日论文日期：2026-05-13；点赞数：19）
  \begin{itemize}[leftmargin=1.5em]
    \item \textbf{动机}：让模型从一次性问答转向长期任务执行，能够主动检索、保存记忆、组合工具并生成可复用技能。
    \item \textbf{背景}：传统检索增强生成与单轮工具调用难以支撑开放任务，长期记忆、隐私、轨迹质量和科研过程建模成为新瓶颈。
    \item \textbf{方法}：论文通常通过高质量轨迹、工具编排、长期记忆结构、多智能体协作或科研图谱，增强智能体的任务闭环。本文题目中的核心对象是“迈向视觉原生多模态深度搜索智能体的在线策略数据演化”，可把它理解为该方向的一种具体实现或问题重述。
    \item \textbf{洞察}：搜索能力的关键不只是召回相似文本，而是智能体如何与语料、工具和反馈形成闭环。
    \item \textbf{结论}：它的结论价值在于把智能体能力从提示词技巧推进到可训练、可评估的系统流程。
  \end{itemize}
  \item \textbf{\href{https://huggingface.co/papers/2605.10344}{TMAS：通过多智能体协同扩展测试时计算}}（\href{https://arxiv.org/abs/2605.10344}{2605.10344}；每日论文日期：2026-05-12；点赞数：46）
  \begin{itemize}[leftmargin=1.5em]
    \item \textbf{动机}：让模型从一次性问答转向长期任务执行，能够主动检索、保存记忆、组合工具并生成可复用技能。
    \item \textbf{背景}：传统检索增强生成与单轮工具调用难以支撑开放任务，长期记忆、隐私、轨迹质量和科研过程建模成为新瓶颈。
    \item \textbf{方法}：论文通常通过高质量轨迹、工具编排、长期记忆结构、多智能体协作或科研图谱，增强智能体的任务闭环。本文题目中的核心对象是“TMAS：通过多智能体协同扩展测试时计算”，可把它理解为该方向的一种具体实现或问题重述。
    \item \textbf{洞察}：这篇工作的关键不在单点技巧，而在它把数据、模型、推理、奖励与评测中的某个环节重新组织进系统闭环。
    \item \textbf{结论}：它的结论价值在于把智能体能力从提示词技巧推进到可训练、可评估的系统流程。
  \end{itemize}
  \item \textbf{\href{https://huggingface.co/papers/2605.09608}{几何冲突：解释并控制大模型持续后训练中的遗忘}}（\href{https://arxiv.org/abs/2605.09608}{2605.09608}；每日论文日期：2026-05-12；点赞数：42）
  \begin{itemize}[leftmargin=1.5em]
    \item \textbf{动机}：让模型从一次性问答转向长期任务执行，能够主动检索、保存记忆、组合工具并生成可复用技能。
    \item \textbf{背景}：传统检索增强生成与单轮工具调用难以支撑开放任务，长期记忆、隐私、轨迹质量和科研过程建模成为新瓶颈。
    \item \textbf{方法}：论文通常通过高质量轨迹、工具编排、长期记忆结构、多智能体协作或科研图谱，增强智能体的任务闭环。本文题目中的核心对象是“几何冲突：解释并控制大模型持续后训练中的遗忘”，可把它理解为该方向的一种具体实现或问题重述。
    \item \textbf{洞察}：这篇工作的关键不在单点技巧，而在它把数据、模型、推理、奖励与评测中的某个环节重新组织进系统闭环。
    \item \textbf{结论}：它的结论价值在于把智能体能力从提示词技巧推进到可训练、可评估的系统流程。
  \end{itemize}
  \item \textbf{\href{https://huggingface.co/papers/2605.08083}{大模型改进大模型：面向测试时扩展的智能体式发现}}（\href{https://arxiv.org/abs/2605.08083}{2605.08083}；每日论文日期：2026-05-11；点赞数：62）
  \begin{itemize}[leftmargin=1.5em]
    \item \textbf{动机}：让模型从一次性问答转向长期任务执行，能够主动检索、保存记忆、组合工具并生成可复用技能。
    \item \textbf{背景}：传统检索增强生成与单轮工具调用难以支撑开放任务，长期记忆、隐私、轨迹质量和科研过程建模成为新瓶颈。
    \item \textbf{方法}：论文通常通过高质量轨迹、工具编排、长期记忆结构、多智能体协作或科研图谱，增强智能体的任务闭环。本文题目中的核心对象是“大模型改进大模型：面向测试时扩展的智能体式发现”，可把它理解为该方向的一种具体实现或问题重述。
    \item \textbf{洞察}：这篇工作的关键不在单点技巧，而在它把数据、模型、推理、奖励与评测中的某个环节重新组织进系统闭环。
    \item \textbf{结论}：它的结论价值在于把智能体能力从提示词技巧推进到可训练、可评估的系统流程。
  \end{itemize}
  \item \textbf{\href{https://huggingface.co/papers/2605.07177}{HyperEyes：面向并行多模态搜索智能体的双粒度效率感知强化学习}}（\href{https://arxiv.org/abs/2605.07177}{2605.07177}；每日论文日期：2026-05-11；点赞数：57）
  \begin{itemize}[leftmargin=1.5em]
    \item \textbf{动机}：让模型从一次性问答转向长期任务执行，能够主动检索、保存记忆、组合工具并生成可复用技能。
    \item \textbf{背景}：传统检索增强生成与单轮工具调用难以支撑开放任务，长期记忆、隐私、轨迹质量和科研过程建模成为新瓶颈。
    \item \textbf{方法}：论文通常通过高质量轨迹、工具编排、长期记忆结构、多智能体协作或科研图谱，增强智能体的任务闭环。本文题目中的核心对象是“HyperEyes：面向并行多模态搜索智能体的双粒度效率感知强化学习”，可把它理解为该方向的一种具体实现或问题重述。
    \item \textbf{洞察}：搜索能力的关键不只是召回相似文本，而是智能体如何与语料、工具和反馈形成闭环。
    \item \textbf{结论}：它的结论价值在于把智能体能力从提示词技巧推进到可训练、可评估的系统流程。
  \end{itemize}
  \item \textbf{\href{https://huggingface.co/papers/2605.04615}{超越检索：面向代码搜索的多任务基准与模型}}（\href{https://arxiv.org/abs/2605.04615}{2605.04615}；每日论文日期：2026-05-11；点赞数：22）
  \begin{itemize}[leftmargin=1.5em]
    \item \textbf{动机}：让模型从一次性问答转向长期任务执行，能够主动检索、保存记忆、组合工具并生成可复用技能。
    \item \textbf{背景}：传统检索增强生成与单轮工具调用难以支撑开放任务，长期记忆、隐私、轨迹质量和科研过程建模成为新瓶颈。
    \item \textbf{方法}：论文通常通过高质量轨迹、工具编排、长期记忆结构、多智能体协作或科研图谱，增强智能体的任务闭环。本文题目中的核心对象是“超越检索：面向代码搜索的多任务基准与模型”，可把它理解为该方向的一种具体实现或问题重述。
    \item \textbf{洞察}：搜索能力的关键不只是召回相似文本，而是智能体如何与语料、工具和反馈形成闭环。
    \item \textbf{结论}：它的结论价值在于把智能体能力从提示词技巧推进到可训练、可评估的系统流程。
  \end{itemize}
  \item \textbf{\href{https://huggingface.co/papers/2605.05242}{超越语义相似度：通过直接语料交互重新思考智能体搜索中的检索}}（\href{https://arxiv.org/abs/2605.05242}{2605.05242}；每日论文日期：2026-05-08；点赞数：101）
  \begin{itemize}[leftmargin=1.5em]
    \item \textbf{动机}：让模型从一次性问答转向长期任务执行，能够主动检索、保存记忆、组合工具并生成可复用技能。
    \item \textbf{背景}：传统检索增强生成与单轮工具调用难以支撑开放任务，长期记忆、隐私、轨迹质量和科研过程建模成为新瓶颈。
    \item \textbf{方法}：论文通常通过高质量轨迹、工具编排、长期记忆结构、多智能体协作或科研图谱，增强智能体的任务闭环。本文题目中的核心对象是“超越语义相似度：通过直接语料交互重新思考智能体搜索中的检索”，可把它理解为该方向的一种具体实现或问题重述。
    \item \textbf{洞察}：真正有用的世界模型必须同时回答状态会如何变化、动作是否可行以及失败风险在哪里。
    \item \textbf{结论}：它的结论价值在于把智能体能力从提示词技巧推进到可训练、可评估的系统流程。
  \end{itemize}
  \item \textbf{\href{https://huggingface.co/papers/2605.06130}{Skill1：通过强化学习统一演化技能增强智能体}}（\href{https://arxiv.org/abs/2605.06130}{2605.06130}；每日论文日期：2026-05-08；点赞数：94）
  \begin{itemize}[leftmargin=1.5em]
    \item \textbf{动机}：让模型从一次性问答转向长期任务执行，能够主动检索、保存记忆、组合工具并生成可复用技能。
    \item \textbf{背景}：传统检索增强生成与单轮工具调用难以支撑开放任务，长期记忆、隐私、轨迹质量和科研过程建模成为新瓶颈。
    \item \textbf{方法}：论文通常通过高质量轨迹、工具编排、长期记忆结构、多智能体协作或科研图谱，增强智能体的任务闭环。本文题目中的核心对象是“Skill1：通过强化学习统一演化技能增强智能体”，可把它理解为该方向的一种具体实现或问题重述。
    \item \textbf{洞察}：奖励信号越能贴近任务结构，后训练越可能稳定地产生能力增益。
    \item \textbf{结论}：它的结论价值在于把智能体能力从提示词技巧推进到可训练、可评估的系统流程。
  \end{itemize}
  \item \textbf{\href{https://huggingface.co/papers/2605.06416}{MiA-Signature：面向长上下文理解的全局激活近似}}（\href{https://arxiv.org/abs/2605.06416}{2605.06416}；每日论文日期：2026-05-08；点赞数：53）
  \begin{itemize}[leftmargin=1.5em]
    \item \textbf{动机}：让模型从一次性问答转向长期任务执行，能够主动检索、保存记忆、组合工具并生成可复用技能。
    \item \textbf{背景}：传统检索增强生成与单轮工具调用难以支撑开放任务，长期记忆、隐私、轨迹质量和科研过程建模成为新瓶颈。
    \item \textbf{方法}：论文通常通过高质量轨迹、工具编排、长期记忆结构、多智能体协作或科研图谱，增强智能体的任务闭环。本文题目中的核心对象是“MiA-Signature：面向长上下文理解的全局激活近似”，可把它理解为该方向的一种具体实现或问题重述。
    \item \textbf{洞察}：这篇工作的关键不在单点技巧，而在它把数据、模型、推理、奖励与评测中的某个环节重新组织进系统闭环。
    \item \textbf{结论}：它的结论价值在于把智能体能力从提示词技巧推进到可训练、可评估的系统流程。
  \end{itemize}
  \item \textbf{\href{https://huggingface.co/papers/2605.06614}{SkillOS：面向自我演化智能体的技能策展学习}}（\href{https://arxiv.org/abs/2605.06614}{2605.06614}；每日论文日期：2026-05-08；点赞数：38）
  \begin{itemize}[leftmargin=1.5em]
    \item \textbf{动机}：让模型从一次性问答转向长期任务执行，能够主动检索、保存记忆、组合工具并生成可复用技能。
    \item \textbf{背景}：传统检索增强生成与单轮工具调用难以支撑开放任务，长期记忆、隐私、轨迹质量和科研过程建模成为新瓶颈。
    \item \textbf{方法}：论文通常通过高质量轨迹、工具编排、长期记忆结构、多智能体协作或科研图谱，增强智能体的任务闭环。本文题目中的核心对象是“SkillOS：面向自我演化智能体的技能策展学习”，可把它理解为该方向的一种具体实现或问题重述。
    \item \textbf{洞察}：这篇工作的关键不在单点技巧，而在它把数据、模型、推理、奖励与评测中的某个环节重新组织进系统闭环。
    \item \textbf{结论}：它的结论价值在于把智能体能力从提示词技巧推进到可训练、可评估的系统流程。
  \end{itemize}
  \item \textbf{\href{https://huggingface.co/papers/2605.05185}{OpenSearch-VL：前沿多模态搜索智能体的开放配方}}（\href{https://arxiv.org/abs/2605.05185}{2605.05185}；每日论文日期：2026-05-07；点赞数：96）
  \begin{itemize}[leftmargin=1.5em]
    \item \textbf{动机}：让模型从一次性问答转向长期任务执行，能够主动检索、保存记忆、组合工具并生成可复用技能。
    \item \textbf{背景}：传统检索增强生成与单轮工具调用难以支撑开放任务，长期记忆、隐私、轨迹质量和科研过程建模成为新瓶颈。
    \item \textbf{方法}：论文通常通过高质量轨迹、工具编排、长期记忆结构、多智能体协作或科研图谱，增强智能体的任务闭环。本文题目中的核心对象是“OpenSearch-VL：前沿多模态搜索智能体的开放配方”，可把它理解为该方向的一种具体实现或问题重述。
    \item \textbf{洞察}：搜索能力的关键不只是召回相似文本，而是智能体如何与语料、工具和反馈形成闭环。
    \item \textbf{结论}：它的结论价值在于把智能体能力从提示词技巧推进到可训练、可评估的系统流程。
  \end{itemize}
  \item \textbf{\href{https://huggingface.co/papers/2605.04018}{重新思考推理密集型检索：评估并推进智能体搜索系统中的检索器}}（\href{https://arxiv.org/abs/2605.04018}{2605.04018}；每日论文日期：2026-05-07；点赞数：38）
  \begin{itemize}[leftmargin=1.5em]
    \item \textbf{动机}：让模型从一次性问答转向长期任务执行，能够主动检索、保存记忆、组合工具并生成可复用技能。
    \item \textbf{背景}：传统检索增强生成与单轮工具调用难以支撑开放任务，长期记忆、隐私、轨迹质量和科研过程建模成为新瓶颈。
    \item \textbf{方法}：论文通常通过高质量轨迹、工具编排、长期记忆结构、多智能体协作或科研图谱，增强智能体的任务闭环。本文题目中的核心对象是“重新思考推理密集型检索：评估并推进智能体搜索系统中的检索器”，可把它理解为该方向的一种具体实现或问题重述。
    \item \textbf{洞察}：搜索能力的关键不只是召回相似文本，而是智能体如何与语料、工具和反馈形成闭环。
    \item \textbf{结论}：它的结论价值在于把智能体能力从提示词技巧推进到可训练、可评估的系统流程。
  \end{itemize}
  \item \textbf{\href{https://huggingface.co/papers/2605.03042}{ARIS：通过对抗式多智能体协作实现自主科研}}（\href{https://arxiv.org/abs/2605.03042}{2605.03042}；每日论文日期：2026-05-06；点赞数：110）
  \begin{itemize}[leftmargin=1.5em]
    \item \textbf{动机}：让模型从一次性问答转向长期任务执行，能够主动检索、保存记忆、组合工具并生成可复用技能。
    \item \textbf{背景}：传统检索增强生成与单轮工具调用难以支撑开放任务，长期记忆、隐私、轨迹质量和科研过程建模成为新瓶颈。
    \item \textbf{方法}：论文通常通过高质量轨迹、工具编排、长期记忆结构、多智能体协作或科研图谱，增强智能体的任务闭环。本文题目中的核心对象是“ARIS：通过对抗式多智能体协作实现自主科研”，可把它理解为该方向的一种具体实现或问题重述。
    \item \textbf{洞察}：搜索能力的关键不只是召回相似文本，而是智能体如何与语料、工具和反馈形成闭环。
    \item \textbf{结论}：它的结论价值在于把智能体能力从提示词技巧推进到可训练、可评估的系统流程。
  \end{itemize}
  \item \textbf{\href{https://huggingface.co/papers/2605.04036}{OpenSeeker-v2：用高信息量困难轨迹推进搜索智能体边界}}（\href{https://arxiv.org/abs/2605.04036}{2605.04036}；每日论文日期：2026-05-06；点赞数：64）
  \begin{itemize}[leftmargin=1.5em]
    \item \textbf{动机}：让模型从一次性问答转向长期任务执行，能够主动检索、保存记忆、组合工具并生成可复用技能。
    \item \textbf{背景}：传统检索增强生成与单轮工具调用难以支撑开放任务，长期记忆、隐私、轨迹质量和科研过程建模成为新瓶颈。
    \item \textbf{方法}：论文通常通过高质量轨迹、工具编排、长期记忆结构、多智能体协作或科研图谱，增强智能体的任务闭环。本文题目中的核心对象是“OpenSeeker-v2：用高信息量困难轨迹推进搜索智能体边界”，可把它理解为该方向的一种具体实现或问题重述。
    \item \textbf{洞察}：搜索能力的关键不只是召回相似文本，而是智能体如何与语料、工具和反馈形成闭环。
    \item \textbf{结论}：它的结论价值在于把智能体能力从提示词技巧推进到可训练、可评估的系统流程。
  \end{itemize}
  \item \textbf{\href{https://huggingface.co/papers/2605.02396}{HeavySkill：把重思考作为智能体内在技能}}（\href{https://arxiv.org/abs/2605.02396}{2605.02396}；每日论文日期：2026-05-06；点赞数：22）
  \begin{itemize}[leftmargin=1.5em]
    \item \textbf{动机}：让模型从一次性问答转向长期任务执行，能够主动检索、保存记忆、组合工具并生成可复用技能。
    \item \textbf{背景}：传统检索增强生成与单轮工具调用难以支撑开放任务，长期记忆、隐私、轨迹质量和科研过程建模成为新瓶颈。
    \item \textbf{方法}：论文通常通过高质量轨迹、工具编排、长期记忆结构、多智能体协作或科研图谱，增强智能体的任务闭环。本文题目中的核心对象是“HeavySkill：把重思考作为智能体内在技能”，可把它理解为该方向的一种具体实现或问题重述。
    \item \textbf{洞察}：这篇工作的关键不在单点技巧，而在它把数据、模型、推理、奖励与评测中的某个环节重新组织进系统闭环。
    \item \textbf{结论}：它的结论价值在于把智能体能力从提示词技巧推进到可训练、可评估的系统流程。
  \end{itemize}
  \item \textbf{\href{https://huggingface.co/papers/2604.27660}{从上下文到技能：语言模型能否从上下文中学会技能}}（\href{https://arxiv.org/abs/2604.27660}{2604.27660}；每日论文日期：2026-05-05；点赞数：153）
  \begin{itemize}[leftmargin=1.5em]
    \item \textbf{动机}：让模型从一次性问答转向长期任务执行，能够主动检索、保存记忆、组合工具并生成可复用技能。
    \item \textbf{背景}：传统检索增强生成与单轮工具调用难以支撑开放任务，长期记忆、隐私、轨迹质量和科研过程建模成为新瓶颈。
    \item \textbf{方法}：论文通常通过高质量轨迹、工具编排、长期记忆结构、多智能体协作或科研图谱，增强智能体的任务闭环。本文题目中的核心对象是“从上下文到技能：语言模型能否从上下文中学会技能”，可把它理解为该方向的一种具体实现或问题重述。
    \item \textbf{洞察}：这篇工作的关键不在单点技巧，而在它把数据、模型、推理、奖励与评测中的某个环节重新组织进系统闭环。
    \item \textbf{结论}：它的结论价值在于把智能体能力从提示词技巧推进到可训练、可评估的系统流程。
  \end{itemize}
  \item \textbf{\href{https://huggingface.co/papers/2604.28075}{重复优先于多样性：面向样本高效德语建模的高信号数据过滤}}（\href{https://arxiv.org/abs/2604.28075}{2604.28075}；每日论文日期：2026-05-05；点赞数：19）
  \begin{itemize}[leftmargin=1.5em]
    \item \textbf{动机}：让模型从一次性问答转向长期任务执行，能够主动检索、保存记忆、组合工具并生成可复用技能。
    \item \textbf{背景}：传统检索增强生成与单轮工具调用难以支撑开放任务，长期记忆、隐私、轨迹质量和科研过程建模成为新瓶颈。
    \item \textbf{方法}：论文通常通过高质量轨迹、工具编排、长期记忆结构、多智能体协作或科研图谱，增强智能体的任务闭环。本文题目中的核心对象是“重复优先于多样性：面向样本高效德语建模的高信号数据过滤”，可把它理解为该方向的一种具体实现或问题重述。
    \item \textbf{洞察}：这篇工作的关键不在单点技巧，而在它把数据、模型、推理、奖励与评测中的某个环节重新组织进系统闭环。
    \item \textbf{结论}：它的结论价值在于把智能体能力从提示词技巧推进到可训练、可评估的系统流程。
  \end{itemize}
  \item \textbf{\href{https://huggingface.co/papers/2605.02661}{AcademiClaw：当学生为 AI 智能体设置挑战}}（\href{https://arxiv.org/abs/2605.02661}{2605.02661}；每日论文日期：2026-05-05；点赞数：16）
  \begin{itemize}[leftmargin=1.5em]
    \item \textbf{动机}：让模型从一次性问答转向长期任务执行，能够主动检索、保存记忆、组合工具并生成可复用技能。
    \item \textbf{背景}：传统检索增强生成与单轮工具调用难以支撑开放任务，长期记忆、隐私、轨迹质量和科研过程建模成为新瓶颈。
    \item \textbf{方法}：论文通常通过高质量轨迹、工具编排、长期记忆结构、多智能体协作或科研图谱，增强智能体的任务闭环。本文题目中的核心对象是“AcademiClaw：当学生为 AI 智能体设置挑战”，可把它理解为该方向的一种具体实现或问题重述。
    \item \textbf{洞察}：这篇工作的关键不在单点技巧，而在它把数据、模型、推理、奖励与评测中的某个环节重新组织进系统闭环。
    \item \textbf{结论}：它的结论价值在于把智能体能力从提示词技巧推进到可训练、可评估的系统流程。
  \end{itemize}
  \item \textbf{\href{https://huggingface.co/papers/2604.27221}{Web2BigTable：面向互联网规模信息搜索与抽取的双层多智能体系统}}（\href{https://arxiv.org/abs/2604.27221}{2604.27221}；每日论文日期：2026-05-04；点赞数：38）
  \begin{itemize}[leftmargin=1.5em]
    \item \textbf{动机}：让模型从一次性问答转向长期任务执行，能够主动检索、保存记忆、组合工具并生成可复用技能。
    \item \textbf{背景}：传统检索增强生成与单轮工具调用难以支撑开放任务，长期记忆、隐私、轨迹质量和科研过程建模成为新瓶颈。
    \item \textbf{方法}：论文通常通过高质量轨迹、工具编排、长期记忆结构、多智能体协作或科研图谱，增强智能体的任务闭环。本文题目中的核心对象是“Web2BigTable：面向互联网规模信息搜索与抽取的双层多智能体系统”，可把它理解为该方向的一种具体实现或问题重述。
    \item \textbf{洞察}：真正有用的世界模型必须同时回答状态会如何变化、动作是否可行以及失败风险在哪里。
    \item \textbf{结论}：它的结论价值在于把智能体能力从提示词技巧推进到可训练、可评估的系统流程。
  \end{itemize}
  \item \textbf{\href{https://huggingface.co/papers/2604.24026}{从技能文本到技能结构：面向智能体技能的调度—结构—逻辑表征}}（\href{https://arxiv.org/abs/2604.24026}{2604.24026}；每日论文日期：2026-05-04；点赞数：21）
  \begin{itemize}[leftmargin=1.5em]
    \item \textbf{动机}：让模型从一次性问答转向长期任务执行，能够主动检索、保存记忆、组合工具并生成可复用技能。
    \item \textbf{背景}：传统检索增强生成与单轮工具调用难以支撑开放任务，长期记忆、隐私、轨迹质量和科研过程建模成为新瓶颈。
    \item \textbf{方法}：论文通常通过高质量轨迹、工具编排、长期记忆结构、多智能体协作或科研图谱，增强智能体的任务闭环。本文题目中的核心对象是“从技能文本到技能结构：面向智能体技能的调度—结构—逻辑表征”，可把它理解为该方向的一种具体实现或问题重述。
    \item \textbf{洞察}：这篇工作的关键不在单点技巧，而在它把数据、模型、推理、奖励与评测中的某个环节重新组织进系统闭环。
    \item \textbf{结论}：它的结论价值在于把智能体能力从提示词技巧推进到可训练、可评估的系统流程。
  \end{itemize}
  \item \textbf{\href{https://huggingface.co/papers/2604.28158}{Intern-Atlas：作为 AI 科学家研究基础设施的方法演化图谱}}（\href{https://arxiv.org/abs/2604.28158}{2604.28158}；每日论文日期：2026-05-01；点赞数：47）
  \begin{itemize}[leftmargin=1.5em]
    \item \textbf{动机}：让模型从一次性问答转向长期任务执行，能够主动检索、保存记忆、组合工具并生成可复用技能。
    \item \textbf{背景}：传统检索增强生成与单轮工具调用难以支撑开放任务，长期记忆、隐私、轨迹质量和科研过程建模成为新瓶颈。
    \item \textbf{方法}：论文通常通过高质量轨迹、工具编排、长期记忆结构、多智能体协作或科研图谱，增强智能体的任务闭环。本文题目中的核心对象是“Intern-Atlas：作为 AI 科学家研究基础设施的方法演化图谱”，可把它理解为该方向的一种具体实现或问题重述。
    \item \textbf{洞察}：搜索能力的关键不只是召回相似文本，而是智能体如何与语料、工具和反馈形成闭环。
    \item \textbf{结论}：它的结论价值在于把智能体能力从提示词技巧推进到可训练、可评估的系统流程。
  \end{itemize}
  \item \textbf{\href{https://huggingface.co/papers/2604.28139}{Claw-Eval-Live：面向演化真实工作流的在线智能体基准}}（\href{https://arxiv.org/abs/2604.28139}{2604.28139}；每日论文日期：2026-05-01；点赞数：42）
  \begin{itemize}[leftmargin=1.5em]
    \item \textbf{动机}：让模型从一次性问答转向长期任务执行，能够主动检索、保存记忆、组合工具并生成可复用技能。
    \item \textbf{背景}：传统检索增强生成与单轮工具调用难以支撑开放任务，长期记忆、隐私、轨迹质量和科研过程建模成为新瓶颈。
    \item \textbf{方法}：论文通常通过高质量轨迹、工具编排、长期记忆结构、多智能体协作或科研图谱，增强智能体的任务闭环。本文题目中的核心对象是“Claw-Eval-Live：面向演化真实工作流的在线智能体基准”，可把它理解为该方向的一种具体实现或问题重述。
    \item \textbf{洞察}：真正有用的世界模型必须同时回答状态会如何变化、动作是否可行以及失败风险在哪里。
    \item \textbf{结论}：它的结论价值在于把智能体能力从提示词技巧推进到可训练、可评估的系统流程。
  \end{itemize}
  \item \textbf{\href{https://huggingface.co/papers/2604.24658}{最后一篇人类撰写的论文：智能体原生研究产物}}（\href{https://arxiv.org/abs/2604.24658}{2604.24658}；每日论文日期：2026-05-01；点赞数：20）
  \begin{itemize}[leftmargin=1.5em]
    \item \textbf{动机}：让模型从一次性问答转向长期任务执行，能够主动检索、保存记忆、组合工具并生成可复用技能。
    \item \textbf{背景}：传统检索增强生成与单轮工具调用难以支撑开放任务，长期记忆、隐私、轨迹质量和科研过程建模成为新瓶颈。
    \item \textbf{方法}：论文通常通过高质量轨迹、工具编排、长期记忆结构、多智能体协作或科研图谱，增强智能体的任务闭环。本文题目中的核心对象是“最后一篇人类撰写的论文：智能体原生研究产物”，可把它理解为该方向的一种具体实现或问题重述。
    \item \textbf{洞察}：搜索能力的关键不只是召回相似文本，而是智能体如何与语料、工具和反馈形成闭环。
    \item \textbf{结论}：它的结论价值在于把智能体能力从提示词技巧推进到可训练、可评估的系统流程。
  \end{itemize}
  \item \textbf{\href{https://huggingface.co/papers/2604.28181}{大规模合成计算机：长程生产力仿真}}（\href{https://arxiv.org/abs/2604.28181}{2604.28181}；每日论文日期：2026-05-01；点赞数：18）
  \begin{itemize}[leftmargin=1.5em]
    \item \textbf{动机}：让模型从一次性问答转向长期任务执行，能够主动检索、保存记忆、组合工具并生成可复用技能。
    \item \textbf{背景}：传统检索增强生成与单轮工具调用难以支撑开放任务，长期记忆、隐私、轨迹质量和科研过程建模成为新瓶颈。
    \item \textbf{方法}：论文通常通过高质量轨迹、工具编排、长期记忆结构、多智能体协作或科研图谱，增强智能体的任务闭环。本文题目中的核心对象是“大规模合成计算机：长程生产力仿真”，可把它理解为该方向的一种具体实现或问题重述。
    \item \textbf{洞察}：这篇工作的关键不在单点技巧，而在它把数据、模型、推理、奖励与评测中的某个环节重新组织进系统闭环。
    \item \textbf{结论}：它的结论价值在于把智能体能力从提示词技巧推进到可训练、可评估的系统流程。
  \end{itemize}
  \item \textbf{\href{https://huggingface.co/papers/2604.27151}{面向高效计算机使用智能体的步骤级优化}}（\href{https://arxiv.org/abs/2604.27151}{2604.27151}；每日论文日期：2026-05-01；点赞数：18）
  \begin{itemize}[leftmargin=1.5em]
    \item \textbf{动机}：让模型从一次性问答转向长期任务执行，能够主动检索、保存记忆、组合工具并生成可复用技能。
    \item \textbf{背景}：传统检索增强生成与单轮工具调用难以支撑开放任务，长期记忆、隐私、轨迹质量和科研过程建模成为新瓶颈。
    \item \textbf{方法}：论文通常通过高质量轨迹、工具编排、长期记忆结构、多智能体协作或科研图谱，增强智能体的任务闭环。本文题目中的核心对象是“面向高效计算机使用智能体的步骤级优化”，可把它理解为该方向的一种具体实现或问题重述。
    \item \textbf{洞察}：这篇工作的关键不在单点技巧，而在它把数据、模型、推理、奖励与评测中的某个环节重新组织进系统闭环。
    \item \textbf{结论}：它的结论价值在于把智能体能力从提示词技巧推进到可训练、可评估的系统流程。
  \end{itemize}
\end{enumerate}

\subsubsection{类别趋势}
趋势：智能体搜索正在从检索增强生成走向直接语料交互、高难轨迹构造、长期记忆压缩、多智能体协作和科研方法图谱。

\subsubsection{下一阶段预测}
预测：未来半年会形成智能体数据飞轮：搜索失败被转化为训练轨迹，记忆隐私成为默认约束，科研智能体会从摘要工具变成假设生成和对抗审稿基础设施。

\subsection{后训练、可验证奖励强化学习、奖励建模与测试时计算}

本类关注奖励、策略优化、在线策略数据演化、评价细则和测试时计算如何改变模型能力。

\subsubsection{论文卡片}
\begin{enumerate}[leftmargin=1.5em]
  \item \textbf{\href{https://huggingface.co/papers/2605.12178}{企业系统是否需要学习型世界模型：上下文对推断动态的重要性}}（\href{https://arxiv.org/abs/2605.12178}{2605.12178}；每日论文日期：2026-05-13；点赞数：53）
  \begin{itemize}[leftmargin=1.5em]
    \item \textbf{动机}：让训练和推理中的反馈信号更细、更准、更可验证，从而稳定提升复杂任务能力。
    \item \textbf{背景}：监督微调与单一偏好分数不足以解释复杂任务表现，研究社区开始把奖励拆成步骤、局部区域、评价细则和测试时策略。
    \item \textbf{方法}：论文通常通过在线策略蒸馏、组级优化、评价细则分解、验证器奖励或测试时搜索，把反馈信号变成可训练的改进方向。本文题目中的核心对象是“企业系统是否需要学习型世界模型：上下文对推断动态的重要性”，可把它理解为该方向的一种具体实现或问题重述。
    \item \textbf{洞察}：真正有用的世界模型必须同时回答状态会如何变化、动作是否可行以及失败风险在哪里。
    \item \textbf{结论}：它的结论价值在于说明后训练收益越来越依赖奖励结构和数据演化，而不是只依赖更大的模型。
  \end{itemize}
  \item \textbf{\href{https://huggingface.co/papers/2605.12495}{AlphaGRPO：通过分解式可验证奖励激发统一多模态模型的自反思生成}}（\href{https://arxiv.org/abs/2605.12495}{2605.12495}；每日论文日期：2026-05-13；点赞数：28）
  \begin{itemize}[leftmargin=1.5em]
    \item \textbf{动机}：让训练和推理中的反馈信号更细、更准、更可验证，从而稳定提升复杂任务能力。
    \item \textbf{背景}：监督微调与单一偏好分数不足以解释复杂任务表现，研究社区开始把奖励拆成步骤、局部区域、评价细则和测试时策略。
    \item \textbf{方法}：论文通常通过在线策略蒸馏、组级优化、评价细则分解、验证器奖励或测试时搜索，把反馈信号变成可训练的改进方向。本文题目中的核心对象是“AlphaGRPO：通过分解式可验证奖励激发统一多模态模型的自反思生成”，可把它理解为该方向的一种具体实现或问题重述。
    \item \textbf{洞察}：奖励信号越能贴近任务结构，后训练越可能稳定地产生能力增益。
    \item \textbf{结论}：它的结论价值在于说明后训练收益越来越依赖奖励结构和数据演化，而不是只依赖更大的模型。
  \end{itemize}
  \item \textbf{\href{https://huggingface.co/papers/2605.07237}{教语言模型用代码思考}}（\href{https://arxiv.org/abs/2605.07237}{2605.07237}；每日论文日期：2026-05-13；点赞数：17）
  \begin{itemize}[leftmargin=1.5em]
    \item \textbf{动机}：让训练和推理中的反馈信号更细、更准、更可验证，从而稳定提升复杂任务能力。
    \item \textbf{背景}：监督微调与单一偏好分数不足以解释复杂任务表现，研究社区开始把奖励拆成步骤、局部区域、评价细则和测试时策略。
    \item \textbf{方法}：论文通常通过在线策略蒸馏、组级优化、评价细则分解、验证器奖励或测试时搜索，把反馈信号变成可训练的改进方向。本文题目中的核心对象是“教语言模型用代码思考”，可把它理解为该方向的一种具体实现或问题重述。
    \item \textbf{洞察}：这篇工作的关键不在单点技巧，而在它把数据、模型、推理、奖励与评测中的某个环节重新组织进系统闭环。
    \item \textbf{结论}：它的结论价值在于说明后训练收益越来越依赖奖励结构和数据演化，而不是只依赖更大的模型。
  \end{itemize}
  \item \textbf{\href{https://huggingface.co/papers/2509.24244}{大语言模型中的模型合并缩放律}}（\href{https://arxiv.org/abs/2509.24244}{2509.24244}；每日论文日期：2026-05-12；点赞数：39）
  \begin{itemize}[leftmargin=1.5em]
    \item \textbf{动机}：让训练和推理中的反馈信号更细、更准、更可验证，从而稳定提升复杂任务能力。
    \item \textbf{背景}：监督微调与单一偏好分数不足以解释复杂任务表现，研究社区开始把奖励拆成步骤、局部区域、评价细则和测试时策略。
    \item \textbf{方法}：论文通常通过在线策略蒸馏、组级优化、评价细则分解、验证器奖励或测试时搜索，把反馈信号变成可训练的改进方向。本文题目中的核心对象是“大语言模型中的模型合并缩放律”，可把它理解为该方向的一种具体实现或问题重述。
    \item \textbf{洞察}：这篇工作的关键不在单点技巧，而在它把数据、模型、推理、奖励与评测中的某个环节重新组织进系统闭环。
    \item \textbf{结论}：它的结论价值在于说明后训练收益越来越依赖奖励结构和数据演化，而不是只依赖更大的模型。
  \end{itemize}
  \item \textbf{\href{https://huggingface.co/papers/2605.07465}{SEIF：面向指令遵循的自演化强化学习}}（\href{https://arxiv.org/abs/2605.07465}{2605.07465}；每日论文日期：2026-05-12；点赞数：25）
  \begin{itemize}[leftmargin=1.5em]
    \item \textbf{动机}：让训练和推理中的反馈信号更细、更准、更可验证，从而稳定提升复杂任务能力。
    \item \textbf{背景}：监督微调与单一偏好分数不足以解释复杂任务表现，研究社区开始把奖励拆成步骤、局部区域、评价细则和测试时策略。
    \item \textbf{方法}：论文通常通过在线策略蒸馏、组级优化、评价细则分解、验证器奖励或测试时搜索，把反馈信号变成可训练的改进方向。本文题目中的核心对象是“SEIF：面向指令遵循的自演化强化学习”，可把它理解为该方向的一种具体实现或问题重述。
    \item \textbf{洞察}：奖励信号越能贴近任务结构，后训练越可能稳定地产生能力增益。
    \item \textbf{结论}：它的结论价值在于说明后训练收益越来越依赖奖励结构和数据演化，而不是只依赖更大的模型。
  \end{itemize}
  \item \textbf{\href{https://huggingface.co/papers/2605.08354}{把自动评价细则作为奖励：从隐式偏好到显式多模态生成标准}}（\href{https://arxiv.org/abs/2605.08354}{2605.08354}；每日论文日期：2026-05-12；点赞数：21）
  \begin{itemize}[leftmargin=1.5em]
    \item \textbf{动机}：让训练和推理中的反馈信号更细、更准、更可验证，从而稳定提升复杂任务能力。
    \item \textbf{背景}：监督微调与单一偏好分数不足以解释复杂任务表现，研究社区开始把奖励拆成步骤、局部区域、评价细则和测试时策略。
    \item \textbf{方法}：论文通常通过在线策略蒸馏、组级优化、评价细则分解、验证器奖励或测试时搜索，把反馈信号变成可训练的改进方向。本文题目中的核心对象是“把自动评价细则作为奖励：从隐式偏好到显式多模态生成标准”，可把它理解为该方向的一种具体实现或问题重述。
    \item \textbf{洞察}：奖励信号越能贴近任务结构，后训练越可能稳定地产生能力增益。
    \item \textbf{结论}：它的结论价值在于说明后训练收益越来越依赖奖励结构和数据演化，而不是只依赖更大的模型。
  \end{itemize}
  \item \textbf{\href{https://huggingface.co/papers/2605.09959}{G-Zero：从零数据开始的开放式生成自博弈}}（\href{https://arxiv.org/abs/2605.09959}{2605.09959}；每日论文日期：2026-05-12；点赞数：15）
  \begin{itemize}[leftmargin=1.5em]
    \item \textbf{动机}：让训练和推理中的反馈信号更细、更准、更可验证，从而稳定提升复杂任务能力。
    \item \textbf{背景}：监督微调与单一偏好分数不足以解释复杂任务表现，研究社区开始把奖励拆成步骤、局部区域、评价细则和测试时策略。
    \item \textbf{方法}：论文通常通过在线策略蒸馏、组级优化、评价细则分解、验证器奖励或测试时搜索，把反馈信号变成可训练的改进方向。本文题目中的核心对象是“G-Zero：从零数据开始的开放式生成自博弈”，可把它理解为该方向的一种具体实现或问题重述。
    \item \textbf{洞察}：生成模型的核心竞争正从单帧质量转向可控性、一致性和工作流可用性。
    \item \textbf{结论}：它的结论价值在于说明后训练收益越来越依赖奖励结构和数据演化，而不是只依赖更大的模型。
  \end{itemize}
  \item \textbf{\href{https://huggingface.co/papers/2605.08063}{Flow-OPD：面向流匹配模型的在线策略蒸馏}}（\href{https://arxiv.org/abs/2605.08063}{2605.08063}；每日论文日期：2026-05-11；点赞数：88）
  \begin{itemize}[leftmargin=1.5em]
    \item \textbf{动机}：让训练和推理中的反馈信号更细、更准、更可验证，从而稳定提升复杂任务能力。
    \item \textbf{背景}：监督微调与单一偏好分数不足以解释复杂任务表现，研究社区开始把奖励拆成步骤、局部区域、评价细则和测试时策略。
    \item \textbf{方法}：论文通常通过在线策略蒸馏、组级优化、评价细则分解、验证器奖励或测试时搜索，把反馈信号变成可训练的改进方向。本文题目中的核心对象是“Flow-OPD：面向流匹配模型的在线策略蒸馏”，可把它理解为该方向的一种具体实现或问题重述。
    \item \textbf{洞察}：奖励信号越能贴近任务结构，后训练越可能稳定地产生能力增益。
    \item \textbf{结论}：它的结论价值在于说明后训练收益越来越依赖奖励结构和数据演化，而不是只依赖更大的模型。
  \end{itemize}
  \item \textbf{\href{https://huggingface.co/papers/2605.06139}{列表式策略优化：把组级可验证奖励强化学习看作回答单纯形上的目标投影}}（\href{https://arxiv.org/abs/2605.06139}{2605.06139}；每日论文日期：2026-05-11；点赞数：63）
  \begin{itemize}[leftmargin=1.5em]
    \item \textbf{动机}：让训练和推理中的反馈信号更细、更准、更可验证，从而稳定提升复杂任务能力。
    \item \textbf{背景}：监督微调与单一偏好分数不足以解释复杂任务表现，研究社区开始把奖励拆成步骤、局部区域、评价细则和测试时策略。
    \item \textbf{方法}：论文通常通过在线策略蒸馏、组级优化、评价细则分解、验证器奖励或测试时搜索，把反馈信号变成可训练的改进方向。本文题目中的核心对象是“列表式策略优化：把组级可验证奖励强化学习看作回答单纯形上的目标投影”，可把它理解为该方向的一种具体实现或问题重述。
    \item \textbf{洞察}：奖励信号越能贴近任务结构，后训练越可能稳定地产生能力增益。
    \item \textbf{结论}：它的结论价值在于说明后训练收益越来越依赖奖励结构和数据演化，而不是只依赖更大的模型。
  \end{itemize}
  \item \textbf{\href{https://huggingface.co/papers/2605.06747}{HumanNet：把以人为中心的视频学习扩展到百万小时}}（\href{https://arxiv.org/abs/2605.06747}{2605.06747}；每日论文日期：2026-05-11；点赞数：48）
  \begin{itemize}[leftmargin=1.5em]
    \item \textbf{动机}：让训练和推理中的反馈信号更细、更准、更可验证，从而稳定提升复杂任务能力。
    \item \textbf{背景}：监督微调与单一偏好分数不足以解释复杂任务表现，研究社区开始把奖励拆成步骤、局部区域、评价细则和测试时策略。
    \item \textbf{方法}：论文通常通过在线策略蒸馏、组级优化、评价细则分解、验证器奖励或测试时搜索，把反馈信号变成可训练的改进方向。本文题目中的核心对象是“HumanNet：把以人为中心的视频学习扩展到百万小时”，可把它理解为该方向的一种具体实现或问题重述。
    \item \textbf{洞察}：生成模型的核心竞争正从单帧质量转向可控性、一致性和工作流可用性。
    \item \textbf{结论}：它的结论价值在于说明后训练收益越来越依赖奖励结构和数据演化，而不是只依赖更大的模型。
  \end{itemize}
  \item \textbf{\href{https://huggingface.co/papers/2605.07396}{基于评价细则的在线策略蒸馏}}（\href{https://arxiv.org/abs/2605.07396}{2605.07396}；每日论文日期：2026-05-11；点赞数：37）
  \begin{itemize}[leftmargin=1.5em]
    \item \textbf{动机}：让训练和推理中的反馈信号更细、更准、更可验证，从而稳定提升复杂任务能力。
    \item \textbf{背景}：监督微调与单一偏好分数不足以解释复杂任务表现，研究社区开始把奖励拆成步骤、局部区域、评价细则和测试时策略。
    \item \textbf{方法}：论文通常通过在线策略蒸馏、组级优化、评价细则分解、验证器奖励或测试时搜索，把反馈信号变成可训练的改进方向。本文题目中的核心对象是“基于评价细则的在线策略蒸馏”，可把它理解为该方向的一种具体实现或问题重述。
    \item \textbf{洞察}：奖励信号越能贴近任务结构，后训练越可能稳定地产生能力增益。
    \item \textbf{结论}：它的结论价值在于说明后训练收益越来越依赖奖励结构和数据演化，而不是只依赖更大的模型。
  \end{itemize}
  \item \textbf{\href{https://huggingface.co/papers/2605.00425}{AEM：面向多轮智能体强化学习的自适应熵调制}}（\href{https://arxiv.org/abs/2605.00425}{2605.00425}；每日论文日期：2026-05-11；点赞数：20）
  \begin{itemize}[leftmargin=1.5em]
    \item \textbf{动机}：让训练和推理中的反馈信号更细、更准、更可验证，从而稳定提升复杂任务能力。
    \item \textbf{背景}：监督微调与单一偏好分数不足以解释复杂任务表现，研究社区开始把奖励拆成步骤、局部区域、评价细则和测试时策略。
    \item \textbf{方法}：论文通常通过在线策略蒸馏、组级优化、评价细则分解、验证器奖励或测试时搜索，把反馈信号变成可训练的改进方向。本文题目中的核心对象是“AEM：面向多轮智能体强化学习的自适应熵调制”，可把它理解为该方向的一种具体实现或问题重述。
    \item \textbf{洞察}：奖励信号越能贴近任务结构，后训练越可能稳定地产生能力增益。
    \item \textbf{结论}：它的结论价值在于说明后训练收益越来越依赖奖励结构和数据演化，而不是只依赖更大的模型。
  \end{itemize}
  \item \textbf{\href{https://huggingface.co/papers/2605.06507}{MARBLE：面向扩散强化学习的多方面奖励平衡}}（\href{https://arxiv.org/abs/2605.06507}{2605.06507}；每日论文日期：2026-05-08；点赞数：37）
  \begin{itemize}[leftmargin=1.5em]
    \item \textbf{动机}：让训练和推理中的反馈信号更细、更准、更可验证，从而稳定提升复杂任务能力。
    \item \textbf{背景}：监督微调与单一偏好分数不足以解释复杂任务表现，研究社区开始把奖励拆成步骤、局部区域、评价细则和测试时策略。
    \item \textbf{方法}：论文通常通过在线策略蒸馏、组级优化、评价细则分解、验证器奖励或测试时搜索，把反馈信号变成可训练的改进方向。本文题目中的核心对象是“MARBLE：面向扩散强化学习的多方面奖励平衡”，可把它理解为该方向的一种具体实现或问题重述。
    \item \textbf{洞察}：奖励信号越能贴近任务结构，后训练越可能稳定地产生能力增益。
    \item \textbf{结论}：它的结论价值在于说明后训练收益越来越依赖奖励结构和数据演化，而不是只依赖更大的模型。
  \end{itemize}
  \item \textbf{\href{https://huggingface.co/papers/2605.05566}{无意义也有帮助：提示空间扰动扩展推理探索}}（\href{https://arxiv.org/abs/2605.05566}{2605.05566}；每日论文日期：2026-05-08；点赞数：35）
  \begin{itemize}[leftmargin=1.5em]
    \item \textbf{动机}：让训练和推理中的反馈信号更细、更准、更可验证，从而稳定提升复杂任务能力。
    \item \textbf{背景}：监督微调与单一偏好分数不足以解释复杂任务表现，研究社区开始把奖励拆成步骤、局部区域、评价细则和测试时策略。
    \item \textbf{方法}：论文通常通过在线策略蒸馏、组级优化、评价细则分解、验证器奖励或测试时搜索，把反馈信号变成可训练的改进方向。本文题目中的核心对象是“无意义也有帮助：提示空间扰动扩展推理探索”，可把它理解为该方向的一种具体实现或问题重述。
    \item \textbf{洞察}：这篇工作的关键不在单点技巧，而在它把数据、模型、推理、奖励与评测中的某个环节重新组织进系统闭环。
    \item \textbf{结论}：它的结论价值在于说明后训练收益越来越依赖奖励结构和数据演化，而不是只依赖更大的模型。
  \end{itemize}
  \item \textbf{\href{https://huggingface.co/papers/2605.06376}{少步扩散蒸馏的连续时间分布匹配}}（\href{https://arxiv.org/abs/2605.06376}{2605.06376}；每日论文日期：2026-05-08；点赞数：25）
  \begin{itemize}[leftmargin=1.5em]
    \item \textbf{动机}：让训练和推理中的反馈信号更细、更准、更可验证，从而稳定提升复杂任务能力。
    \item \textbf{背景}：监督微调与单一偏好分数不足以解释复杂任务表现，研究社区开始把奖励拆成步骤、局部区域、评价细则和测试时策略。
    \item \textbf{方法}：论文通常通过在线策略蒸馏、组级优化、评价细则分解、验证器奖励或测试时搜索，把反馈信号变成可训练的改进方向。本文题目中的核心对象是“少步扩散蒸馏的连续时间分布匹配”，可把它理解为该方向的一种具体实现或问题重述。
    \item \textbf{洞察}：奖励信号越能贴近任务结构，后训练越可能稳定地产生能力增益。
    \item \textbf{结论}：它的结论价值在于说明后训练收益越来越依赖奖励结构和数据演化，而不是只依赖更大的模型。
  \end{itemize}
  \item \textbf{\href{https://huggingface.co/papers/2605.06642}{StraTA：用战略轨迹抽象激励智能体强化学习}}（\href{https://arxiv.org/abs/2605.06642}{2605.06642}；每日论文日期：2026-05-08；点赞数：24）
  \begin{itemize}[leftmargin=1.5em]
    \item \textbf{动机}：让训练和推理中的反馈信号更细、更准、更可验证，从而稳定提升复杂任务能力。
    \item \textbf{背景}：监督微调与单一偏好分数不足以解释复杂任务表现，研究社区开始把奖励拆成步骤、局部区域、评价细则和测试时策略。
    \item \textbf{方法}：论文通常通过在线策略蒸馏、组级优化、评价细则分解、验证器奖励或测试时搜索，把反馈信号变成可训练的改进方向。本文题目中的核心对象是“StraTA：用战略轨迹抽象激励智能体强化学习”，可把它理解为该方向的一种具体实现或问题重述。
    \item \textbf{洞察}：真正有用的世界模型必须同时回答状态会如何变化、动作是否可行以及失败风险在哪里。
    \item \textbf{结论}：它的结论价值在于说明后训练收益越来越依赖奖励结构和数据演化，而不是只依赖更大的模型。
  \end{itemize}
  \item \textbf{\href{https://huggingface.co/papers/2605.03849}{Stream-R1：面向流式视频生成的可靠性—困惑度感知奖励蒸馏}}（\href{https://arxiv.org/abs/2605.03849}{2605.03849}；每日论文日期：2026-05-07；点赞数：122）
  \begin{itemize}[leftmargin=1.5em]
    \item \textbf{动机}：让训练和推理中的反馈信号更细、更准、更可验证，从而稳定提升复杂任务能力。
    \item \textbf{背景}：监督微调与单一偏好分数不足以解释复杂任务表现，研究社区开始把奖励拆成步骤、局部区域、评价细则和测试时策略。
    \item \textbf{方法}：论文通常通过在线策略蒸馏、组级优化、评价细则分解、验证器奖励或测试时搜索，把反馈信号变成可训练的改进方向。本文题目中的核心对象是“Stream-R1：面向流式视频生成的可靠性—困惑度感知奖励蒸馏”，可把它理解为该方向的一种具体实现或问题重述。
    \item \textbf{洞察}：奖励信号越能贴近任务结构，后训练越可能稳定地产生能力增益。
    \item \textbf{结论}：它的结论价值在于说明后训练收益越来越依赖奖励结构和数据演化，而不是只依赖更大的模型。
  \end{itemize}
  \item \textbf{\href{https://huggingface.co/papers/2605.05204}{D-OPSD：面向持续调优步蒸馏扩散模型的在线自蒸馏}}（\href{https://arxiv.org/abs/2605.05204}{2605.05204}；每日论文日期：2026-05-07；点赞数：25）
  \begin{itemize}[leftmargin=1.5em]
    \item \textbf{动机}：让训练和推理中的反馈信号更细、更准、更可验证，从而稳定提升复杂任务能力。
    \item \textbf{背景}：监督微调与单一偏好分数不足以解释复杂任务表现，研究社区开始把奖励拆成步骤、局部区域、评价细则和测试时策略。
    \item \textbf{方法}：论文通常通过在线策略蒸馏、组级优化、评价细则分解、验证器奖励或测试时搜索，把反馈信号变成可训练的改进方向。本文题目中的核心对象是“D-OPSD：面向持续调优步蒸馏扩散模型的在线自蒸馏”，可把它理解为该方向的一种具体实现或问题重述。
    \item \textbf{洞察}：奖励信号越能贴近任务结构，后训练越可能稳定地产生能力增益。
    \item \textbf{结论}：它的结论价值在于说明后训练收益越来越依赖奖励结构和数据演化，而不是只依赖更大的模型。
  \end{itemize}
  \item \textbf{\href{https://huggingface.co/papers/2604.28123}{超越监督微调到强化学习：面向多模态强化学习的黑盒在线策略蒸馏预对齐}}（\href{https://arxiv.org/abs/2604.28123}{2604.28123}；每日论文日期：2026-05-06；点赞数：47）
  \begin{itemize}[leftmargin=1.5em]
    \item \textbf{动机}：让训练和推理中的反馈信号更细、更准、更可验证，从而稳定提升复杂任务能力。
    \item \textbf{背景}：监督微调与单一偏好分数不足以解释复杂任务表现，研究社区开始把奖励拆成步骤、局部区域、评价细则和测试时策略。
    \item \textbf{方法}：论文通常通过在线策略蒸馏、组级优化、评价细则分解、验证器奖励或测试时搜索，把反馈信号变成可训练的改进方向。本文题目中的核心对象是“超越监督微调到强化学习：面向多模态强化学习的黑盒在线策略蒸馏预对齐”，可把它理解为该方向的一种具体实现或问题重述。
    \item \textbf{洞察}：奖励信号越能贴近任务结构，后训练越可能稳定地产生能力增益。
    \item \textbf{结论}：它的结论价值在于说明后训练收益越来越依赖奖励结构和数据演化，而不是只依赖更大的模型。
  \end{itemize}
  \item \textbf{\href{https://huggingface.co/papers/2605.00347}{Odysseus：通过强化学习把视觉语言模型扩展到百轮以上游戏决策}}（\href{https://arxiv.org/abs/2605.00347}{2605.00347}；每日论文日期：2026-05-04；点赞数：16）
  \begin{itemize}[leftmargin=1.5em]
    \item \textbf{动机}：让训练和推理中的反馈信号更细、更准、更可验证，从而稳定提升复杂任务能力。
    \item \textbf{背景}：监督微调与单一偏好分数不足以解释复杂任务表现，研究社区开始把奖励拆成步骤、局部区域、评价细则和测试时策略。
    \item \textbf{方法}：论文通常通过在线策略蒸馏、组级优化、评价细则分解、验证器奖励或测试时搜索，把反馈信号变成可训练的改进方向。本文题目中的核心对象是“Odysseus：通过强化学习把视觉语言模型扩展到百轮以上游戏决策”，可把它理解为该方向的一种具体实现或问题重述。
    \item \textbf{洞察}：奖励信号越能贴近任务结构，后训练越可能稳定地产生能力增益。
    \item \textbf{结论}：它的结论价值在于说明后训练收益越来越依赖奖励结构和数据演化，而不是只依赖更大的模型。
  \end{itemize}
  \item \textbf{\href{https://huggingface.co/papers/2604.27351}{异构科学基础模型协作}}（\href{https://arxiv.org/abs/2604.27351}{2604.27351}；每日论文日期：2026-05-01；点赞数：212）
  \begin{itemize}[leftmargin=1.5em]
    \item \textbf{动机}：让训练和推理中的反馈信号更细、更准、更可验证，从而稳定提升复杂任务能力。
    \item \textbf{背景}：监督微调与单一偏好分数不足以解释复杂任务表现，研究社区开始把奖励拆成步骤、局部区域、评价细则和测试时策略。
    \item \textbf{方法}：论文通常通过在线策略蒸馏、组级优化、评价细则分解、验证器奖励或测试时搜索，把反馈信号变成可训练的改进方向。本文题目中的核心对象是“异构科学基础模型协作”，可把它理解为该方向的一种具体实现或问题重述。
    \item \textbf{洞察}：这篇工作的关键不在单点技巧，而在它把数据、模型、推理、奖励与评测中的某个环节重新组织进系统闭环。
    \item \textbf{结论}：它的结论价值在于说明后训练收益越来越依赖奖励结构和数据演化，而不是只依赖更大的模型。
  \end{itemize}
  \item \textbf{\href{https://huggingface.co/papers/2604.27083}{协同演化的策略蒸馏}}（\href{https://arxiv.org/abs/2604.27083}{2604.27083}；每日论文日期：2026-05-01；点赞数：64）
  \begin{itemize}[leftmargin=1.5em]
    \item \textbf{动机}：让训练和推理中的反馈信号更细、更准、更可验证，从而稳定提升复杂任务能力。
    \item \textbf{背景}：监督微调与单一偏好分数不足以解释复杂任务表现，研究社区开始把奖励拆成步骤、局部区域、评价细则和测试时策略。
    \item \textbf{方法}：论文通常通过在线策略蒸馏、组级优化、评价细则分解、验证器奖励或测试时搜索，把反馈信号变成可训练的改进方向。本文题目中的核心对象是“协同演化的策略蒸馏”，可把它理解为该方向的一种具体实现或问题重述。
    \item \textbf{洞察}：奖励信号越能贴近任务结构，后训练越可能稳定地产生能力增益。
    \item \textbf{结论}：它的结论价值在于说明后训练收益越来越依赖奖励结构和数据演化，而不是只依赖更大的模型。
  \end{itemize}
  \item \textbf{\href{https://huggingface.co/papers/2604.27505}{在图像编辑中利用基于验证器的强化学习}}（\href{https://arxiv.org/abs/2604.27505}{2604.27505}；每日论文日期：2026-05-01；点赞数：57）
  \begin{itemize}[leftmargin=1.5em]
    \item \textbf{动机}：让训练和推理中的反馈信号更细、更准、更可验证，从而稳定提升复杂任务能力。
    \item \textbf{背景}：监督微调与单一偏好分数不足以解释复杂任务表现，研究社区开始把奖励拆成步骤、局部区域、评价细则和测试时策略。
    \item \textbf{方法}：论文通常通过在线策略蒸馏、组级优化、评价细则分解、验证器奖励或测试时搜索，把反馈信号变成可训练的改进方向。本文题目中的核心对象是“在图像编辑中利用基于验证器的强化学习”，可把它理解为该方向的一种具体实现或问题重述。
    \item \textbf{洞察}：奖励信号越能贴近任务结构，后训练越可能稳定地产生能力增益。
    \item \textbf{结论}：它的结论价值在于说明后训练收益越来越依赖奖励结构和数据演化，而不是只依赖更大的模型。
  \end{itemize}
  \item \textbf{\href{https://huggingface.co/papers/2604.27039}{长度价值模型：面向词元级长度建模的可扩展价值预训练}}（\href{https://arxiv.org/abs/2604.27039}{2604.27039}；每日论文日期：2026-05-01；点赞数：24）
  \begin{itemize}[leftmargin=1.5em]
    \item \textbf{动机}：让训练和推理中的反馈信号更细、更准、更可验证，从而稳定提升复杂任务能力。
    \item \textbf{背景}：监督微调与单一偏好分数不足以解释复杂任务表现，研究社区开始把奖励拆成步骤、局部区域、评价细则和测试时策略。
    \item \textbf{方法}：论文通常通过在线策略蒸馏、组级优化、评价细则分解、验证器奖励或测试时搜索，把反馈信号变成可训练的改进方向。本文题目中的核心对象是“长度价值模型：面向词元级长度建模的可扩展价值预训练”，可把它理解为该方向的一种具体实现或问题重述。
    \item \textbf{洞察}：这篇工作的关键不在单点技巧，而在它把数据、模型、推理、奖励与评测中的某个环节重新组织进系统闭环。
    \item \textbf{结论}：它的结论价值在于说明后训练收益越来越依赖奖励结构和数据演化，而不是只依赖更大的模型。
  \end{itemize}
\end{enumerate}

\subsubsection{类别趋势}
趋势：奖励信号正在分解化、过程化、局部化；可验证奖励强化学习不再局限于数学答案，而进入图像编辑、视频生成、多模态搜索和测试时适配。

\subsubsection{下一阶段预测}
预测：下一阶段最重要的不是优化算法名字变化，而是奖励结构工程：评价细则自动生成、时空信用分配、在线策略数据演化和失败样本主动采样。

\subsection{统一多模态理解生成、图像与视频生成}

本类关注图像、视频、音视觉和统一多模态模型如何从生成质量走向可控编辑与统一交互。

\subsubsection{论文卡片}
\begin{enumerate}[leftmargin=1.5em]
  \item \textbf{\href{https://huggingface.co/papers/2605.13062}{Edit-Compass 与 EditReward-Compass：图像编辑和奖励建模的统一基准}}（\href{https://arxiv.org/abs/2605.13062}{2605.13062}；每日论文日期：2026-05-14；点赞数：18）
  \begin{itemize}[leftmargin=1.5em]
    \item \textbf{动机}：把看图、懂视频、生成、编辑、搜索和交互纳入统一接口，并提升长程一致性与可控性。
    \item \textbf{背景}：多模态模型正在从单一理解或单一生成任务走向统一系统，但跨模态一致性、可控编辑和长视频稳定性仍是难点。
    \item \textbf{方法}：论文通常通过统一架构、扩散先验、奖励驱动编辑、视频数据规模化或跨模态条件机制，提升生成和理解的一致性。本文题目中的核心对象是“Edit-Compass 与 EditReward-Compass：图像编辑和奖励建模的统一基准”，可把它理解为该方向的一种具体实现或问题重述。
    \item \textbf{洞察}：奖励信号越能贴近任务结构，后训练越可能稳定地产生能力增益。
    \item \textbf{结论}：它的结论价值在于把多模态生成从视觉质量竞争推向可控、可交互、可生产的系统能力。
  \end{itemize}
  \item \textbf{\href{https://huggingface.co/papers/2605.12500}{SenseNova-U1：用统一架构贯通多模态理解与生成}}（\href{https://arxiv.org/abs/2605.12500}{2605.12500}；每日论文日期：2026-05-13；点赞数：114）
  \begin{itemize}[leftmargin=1.5em]
    \item \textbf{动机}：把看图、懂视频、生成、编辑、搜索和交互纳入统一接口，并提升长程一致性与可控性。
    \item \textbf{背景}：多模态模型正在从单一理解或单一生成任务走向统一系统，但跨模态一致性、可控编辑和长视频稳定性仍是难点。
    \item \textbf{方法}：论文通常通过统一架构、扩散先验、奖励驱动编辑、视频数据规模化或跨模态条件机制，提升生成和理解的一致性。本文题目中的核心对象是“SenseNova-U1：用统一架构贯通多模态理解与生成”，可把它理解为该方向的一种具体实现或问题重述。
    \item \textbf{洞察}：生成模型的核心竞争正从单帧质量转向可控性、一致性和工作流可用性。
    \item \textbf{结论}：它的结论价值在于把多模态生成从视觉质量竞争推向可控、可交互、可生产的系统能力。
  \end{itemize}
  \item \textbf{\href{https://huggingface.co/papers/2605.12496}{CausalCine：面向多镜头视频叙事的实时自回归生成}}（\href{https://arxiv.org/abs/2605.12496}{2605.12496}；每日论文日期：2026-05-13；点赞数：20）
  \begin{itemize}[leftmargin=1.5em]
    \item \textbf{动机}：把看图、懂视频、生成、编辑、搜索和交互纳入统一接口，并提升长程一致性与可控性。
    \item \textbf{背景}：多模态模型正在从单一理解或单一生成任务走向统一系统，但跨模态一致性、可控编辑和长视频稳定性仍是难点。
    \item \textbf{方法}：论文通常通过统一架构、扩散先验、奖励驱动编辑、视频数据规模化或跨模态条件机制，提升生成和理解的一致性。本文题目中的核心对象是“CausalCine：面向多镜头视频叙事的实时自回归生成”，可把它理解为该方向的一种具体实现或问题重述。
    \item \textbf{洞察}：生成模型的核心竞争正从单帧质量转向可控性、一致性和工作流可用性。
    \item \textbf{结论}：它的结论价值在于把多模态生成从视觉质量竞争推向可控、可交互、可生产的系统能力。
  \end{itemize}
  \item \textbf{\href{https://huggingface.co/papers/2605.10730}{Qwen-Image-2.0 技术报告}}（\href{https://arxiv.org/abs/2605.10730}{2605.10730}；每日论文日期：2026-05-12；点赞数：94）
  \begin{itemize}[leftmargin=1.5em]
    \item \textbf{动机}：把看图、懂视频、生成、编辑、搜索和交互纳入统一接口，并提升长程一致性与可控性。
    \item \textbf{背景}：多模态模型正在从单一理解或单一生成任务走向统一系统，但跨模态一致性、可控编辑和长视频稳定性仍是难点。
    \item \textbf{方法}：论文通常通过统一架构、扩散先验、奖励驱动编辑、视频数据规模化或跨模态条件机制，提升生成和理解的一致性。本文题目中的核心对象是“Qwen-Image-2.0 技术报告”，可把它理解为该方向的一种具体实现或问题重述。
    \item \textbf{洞察}：生成模型的核心竞争正从单帧质量转向可控性、一致性和工作流可用性。
    \item \textbf{结论}：它的结论价值在于把多模态生成从视觉质量竞争推向可控、可交互、可生产的系统能力。
  \end{itemize}
  \item \textbf{\href{https://huggingface.co/papers/2605.08735}{CollabVR：视觉语言模型与视频生成模型协同的视频推理}}（\href{https://arxiv.org/abs/2605.08735}{2605.08735}；每日论文日期：2026-05-12；点赞数：60）
  \begin{itemize}[leftmargin=1.5em]
    \item \textbf{动机}：把看图、懂视频、生成、编辑、搜索和交互纳入统一接口，并提升长程一致性与可控性。
    \item \textbf{背景}：多模态模型正在从单一理解或单一生成任务走向统一系统，但跨模态一致性、可控编辑和长视频稳定性仍是难点。
    \item \textbf{方法}：论文通常通过统一架构、扩散先验、奖励驱动编辑、视频数据规模化或跨模态条件机制，提升生成和理解的一致性。本文题目中的核心对象是“CollabVR：视觉语言模型与视频生成模型协同的视频推理”，可把它理解为该方向的一种具体实现或问题重述。
    \item \textbf{洞察}：生成模型的核心竞争正从单帧质量转向可控性、一致性和工作流可用性。
    \item \textbf{结论}：它的结论价值在于把多模态生成从视觉质量竞争推向可控、可交互、可生产的系统能力。
  \end{itemize}
  \item \textbf{\href{https://huggingface.co/papers/2605.10434}{WorldReasonBench：把视频生成器作为未来世界状态预测器进行人类对齐压力测试}}（\href{https://arxiv.org/abs/2605.10434}{2605.10434}；每日论文日期：2026-05-12；点赞数：26）
  \begin{itemize}[leftmargin=1.5em]
    \item \textbf{动机}：把看图、懂视频、生成、编辑、搜索和交互纳入统一接口，并提升长程一致性与可控性。
    \item \textbf{背景}：多模态模型正在从单一理解或单一生成任务走向统一系统，但跨模态一致性、可控编辑和长视频稳定性仍是难点。
    \item \textbf{方法}：论文通常通过统一架构、扩散先验、奖励驱动编辑、视频数据规模化或跨模态条件机制，提升生成和理解的一致性。本文题目中的核心对象是“WorldReasonBench：把视频生成器作为未来世界状态预测器进行人类对齐压力测试”，可把它理解为该方向的一种具体实现或问题重述。
    \item \textbf{洞察}：真正有用的世界模型必须同时回答状态会如何变化、动作是否可行以及失败风险在哪里。
    \item \textbf{结论}：它的结论价值在于把多模态生成从视觉质量竞争推向可控、可交互、可生产的系统能力。
  \end{itemize}
  \item \textbf{\href{https://huggingface.co/papers/2605.08985}{LLaVA-UHD v4：高效多模态大模型视觉编码的关键因素}}（\href{https://arxiv.org/abs/2605.08985}{2605.08985}；每日论文日期：2026-05-12；点赞数：16）
  \begin{itemize}[leftmargin=1.5em]
    \item \textbf{动机}：把看图、懂视频、生成、编辑、搜索和交互纳入统一接口，并提升长程一致性与可控性。
    \item \textbf{背景}：多模态模型正在从单一理解或单一生成任务走向统一系统，但跨模态一致性、可控编辑和长视频稳定性仍是难点。
    \item \textbf{方法}：论文通常通过统一架构、扩散先验、奖励驱动编辑、视频数据规模化或跨模态条件机制，提升生成和理解的一致性。本文题目中的核心对象是“LLaVA-UHD v4：高效多模态大模型视觉编码的关键因素”，可把它理解为该方向的一种具体实现或问题重述。
    \item \textbf{洞察}：生成模型的核心竞争正从单帧质量转向可控性、一致性和工作流可用性。
    \item \textbf{结论}：它的结论价值在于把多模态生成从视觉质量竞争推向可控、可交互、可生产的系统能力。
  \end{itemize}
  \item \textbf{\href{https://huggingface.co/papers/2512.18181}{MACE-Dance：面向音乐驱动舞蹈视频生成的运动—外观级联专家}}（\href{https://arxiv.org/abs/2512.18181}{2512.18181}；每日论文日期：2026-05-11；点赞数：84）
  \begin{itemize}[leftmargin=1.5em]
    \item \textbf{动机}：把看图、懂视频、生成、编辑、搜索和交互纳入统一接口，并提升长程一致性与可控性。
    \item \textbf{背景}：多模态模型正在从单一理解或单一生成任务走向统一系统，但跨模态一致性、可控编辑和长视频稳定性仍是难点。
    \item \textbf{方法}：论文通常通过统一架构、扩散先验、奖励驱动编辑、视频数据规模化或跨模态条件机制，提升生成和理解的一致性。本文题目中的核心对象是“MACE-Dance：面向音乐驱动舞蹈视频生成的运动—外观级联专家”，可把它理解为该方向的一种具体实现或问题重述。
    \item \textbf{洞察}：生成模型的核心竞争正从单帧质量转向可控性、一致性和工作流可用性。
    \item \textbf{结论}：它的结论价值在于把多模态生成从视觉质量竞争推向可控、可交互、可生产的系统能力。
  \end{itemize}
  \item \textbf{\href{https://huggingface.co/papers/2605.05997}{4DThinker：用 4D 影像进行动态空间理解}}（\href{https://arxiv.org/abs/2605.05997}{2605.05997}；每日论文日期：2026-05-11；点赞数：16）
  \begin{itemize}[leftmargin=1.5em]
    \item \textbf{动机}：把看图、懂视频、生成、编辑、搜索和交互纳入统一接口，并提升长程一致性与可控性。
    \item \textbf{背景}：多模态模型正在从单一理解或单一生成任务走向统一系统，但跨模态一致性、可控编辑和长视频稳定性仍是难点。
    \item \textbf{方法}：论文通常通过统一架构、扩散先验、奖励驱动编辑、视频数据规模化或跨模态条件机制，提升生成和理解的一致性。本文题目中的核心对象是“4DThinker：用 4D 影像进行动态空间理解”，可把它理解为该方向的一种具体实现或问题重述。
    \item \textbf{洞察}：真正有用的世界模型必须同时回答状态会如何变化、动作是否可行以及失败风险在哪里。
    \item \textbf{结论}：它的结论价值在于把多模态生成从视觉质量竞争推向可控、可交互、可生产的系统能力。
  \end{itemize}
  \item \textbf{\href{https://huggingface.co/papers/2605.06924}{A2RD：面向长视频一致性的智能体式自回归扩散}}（\href{https://arxiv.org/abs/2605.06924}{2605.06924}；每日论文日期：2026-05-11；点赞数：15）
  \begin{itemize}[leftmargin=1.5em]
    \item \textbf{动机}：把看图、懂视频、生成、编辑、搜索和交互纳入统一接口，并提升长程一致性与可控性。
    \item \textbf{背景}：多模态模型正在从单一理解或单一生成任务走向统一系统，但跨模态一致性、可控编辑和长视频稳定性仍是难点。
    \item \textbf{方法}：论文通常通过统一架构、扩散先验、奖励驱动编辑、视频数据规模化或跨模态条件机制，提升生成和理解的一致性。本文题目中的核心对象是“A2RD：面向长视频一致性的智能体式自回归扩散”，可把它理解为该方向的一种具体实现或问题重述。
    \item \textbf{洞察}：生成模型的核心竞争正从单帧质量转向可控性、一致性和工作流可用性。
    \item \textbf{结论}：它的结论价值在于把多模态生成从视觉质量竞争推向可控、可交互、可生产的系统能力。
  \end{itemize}
  \item \textbf{\href{https://huggingface.co/papers/2605.04045}{大基础模型中的音视觉智能}}（\href{https://arxiv.org/abs/2605.04045}{2605.04045}；每日论文日期：2026-05-08；点赞数：31）
  \begin{itemize}[leftmargin=1.5em]
    \item \textbf{动机}：把看图、懂视频、生成、编辑、搜索和交互纳入统一接口，并提升长程一致性与可控性。
    \item \textbf{背景}：多模态模型正在从单一理解或单一生成任务走向统一系统，但跨模态一致性、可控编辑和长视频稳定性仍是难点。
    \item \textbf{方法}：论文通常通过统一架构、扩散先验、奖励驱动编辑、视频数据规模化或跨模态条件机制，提升生成和理解的一致性。本文题目中的核心对象是“大基础模型中的音视觉智能”，可把它理解为该方向的一种具体实现或问题重述。
    \item \textbf{洞察}：生成模型的核心竞争正从单帧质量转向可控性、一致性和工作流可用性。
    \item \textbf{结论}：它的结论价值在于把多模态生成从视觉质量竞争推向可控、可交互、可生产的系统能力。
  \end{itemize}
  \item \textbf{\href{https://huggingface.co/papers/2605.04461}{Stream-T1：面向流式视频生成的测试时扩展}}（\href{https://arxiv.org/abs/2605.04461}{2605.04461}；每日论文日期：2026-05-07；点赞数：102）
  \begin{itemize}[leftmargin=1.5em]
    \item \textbf{动机}：把看图、懂视频、生成、编辑、搜索和交互纳入统一接口，并提升长程一致性与可控性。
    \item \textbf{背景}：多模态模型正在从单一理解或单一生成任务走向统一系统，但跨模态一致性、可控编辑和长视频稳定性仍是难点。
    \item \textbf{方法}：论文通常通过统一架构、扩散先验、奖励驱动编辑、视频数据规模化或跨模态条件机制，提升生成和理解的一致性。本文题目中的核心对象是“Stream-T1：面向流式视频生成的测试时扩展”，可把它理解为该方向的一种具体实现或问题重述。
    \item \textbf{洞察}：生成模型的核心竞争正从单帧质量转向可控性、一致性和工作流可用性。
    \item \textbf{结论}：它的结论价值在于把多模态生成从视觉质量竞争推向可控、可交互、可生产的系统能力。
  \end{itemize}
  \item \textbf{\href{https://huggingface.co/papers/2604.27393}{MiniCPM-o 4.5：迈向实时全双工全模态交互}}（\href{https://arxiv.org/abs/2604.27393}{2604.27393}；每日论文日期：2026-05-07；点赞数：66）
  \begin{itemize}[leftmargin=1.5em]
    \item \textbf{动机}：把看图、懂视频、生成、编辑、搜索和交互纳入统一接口，并提升长程一致性与可控性。
    \item \textbf{背景}：多模态模型正在从单一理解或单一生成任务走向统一系统，但跨模态一致性、可控编辑和长视频稳定性仍是难点。
    \item \textbf{方法}：论文通常通过统一架构、扩散先验、奖励驱动编辑、视频数据规模化或跨模态条件机制，提升生成和理解的一致性。本文题目中的核心对象是“MiniCPM-o 4.5：迈向实时全双工全模态交互”，可把它理解为该方向的一种具体实现或问题重述。
    \item \textbf{洞察}：真正有用的世界模型必须同时回答状态会如何变化、动作是否可行以及失败风险在哪里。
    \item \textbf{结论}：它的结论价值在于把多模态生成从视觉质量竞争推向可控、可交互、可生产的系统能力。
  \end{itemize}
  \item \textbf{\href{https://huggingface.co/papers/2605.04128}{在统一多模态理解与生成中唤醒空间智能}}（\href{https://arxiv.org/abs/2605.04128}{2605.04128}；每日论文日期：2026-05-07；点赞数：17）
  \begin{itemize}[leftmargin=1.5em]
    \item \textbf{动机}：把看图、懂视频、生成、编辑、搜索和交互纳入统一接口，并提升长程一致性与可控性。
    \item \textbf{背景}：多模态模型正在从单一理解或单一生成任务走向统一系统，但跨模态一致性、可控编辑和长视频稳定性仍是难点。
    \item \textbf{方法}：论文通常通过统一架构、扩散先验、奖励驱动编辑、视频数据规模化或跨模态条件机制，提升生成和理解的一致性。本文题目中的核心对象是“在统一多模态理解与生成中唤醒空间智能”，可把它理解为该方向的一种具体实现或问题重述。
    \item \textbf{洞察}：生成模型的核心竞争正从单帧质量转向可控性、一致性和工作流可用性。
    \item \textbf{结论}：它的结论价值在于把多模态生成从视觉质量竞争推向可控、可交互、可生产的系统能力。
  \end{itemize}
  \item \textbf{\href{https://huggingface.co/papers/2605.00891}{X2SAM：图像与视频中的任意分割}}（\href{https://arxiv.org/abs/2605.00891}{2605.00891}；每日论文日期：2026-05-06；点赞数：25）
  \begin{itemize}[leftmargin=1.5em]
    \item \textbf{动机}：把看图、懂视频、生成、编辑、搜索和交互纳入统一接口，并提升长程一致性与可控性。
    \item \textbf{背景}：多模态模型正在从单一理解或单一生成任务走向统一系统，但跨模态一致性、可控编辑和长视频稳定性仍是难点。
    \item \textbf{方法}：论文通常通过统一架构、扩散先验、奖励驱动编辑、视频数据规模化或跨模态条件机制，提升生成和理解的一致性。本文题目中的核心对象是“X2SAM：图像与视频中的任意分割”，可把它理解为该方向的一种具体实现或问题重述。
    \item \textbf{洞察}：生成模型的核心竞争正从单帧质量转向可控性、一致性和工作流可用性。
    \item \textbf{结论}：它的结论价值在于把多模态生成从视觉质量竞争推向可控、可交互、可生产的系统能力。
  \end{itemize}
  \item \textbf{\href{https://huggingface.co/papers/2605.02134}{使用预测潜变量的视频生成}}（\href{https://arxiv.org/abs/2605.02134}{2605.02134}；每日论文日期：2026-05-06；点赞数：24）
  \begin{itemize}[leftmargin=1.5em]
    \item \textbf{动机}：把看图、懂视频、生成、编辑、搜索和交互纳入统一接口，并提升长程一致性与可控性。
    \item \textbf{背景}：多模态模型正在从单一理解或单一生成任务走向统一系统，但跨模态一致性、可控编辑和长视频稳定性仍是难点。
    \item \textbf{方法}：论文通常通过统一架构、扩散先验、奖励驱动编辑、视频数据规模化或跨模态条件机制，提升生成和理解的一致性。本文题目中的核心对象是“使用预测潜变量的视频生成”，可把它理解为该方向的一种具体实现或问题重述。
    \item \textbf{洞察}：生成模型的核心竞争正从单帧质量转向可控性、一致性和工作流可用性。
    \item \textbf{结论}：它的结论价值在于把多模态生成从视觉质量竞争推向可控、可交互、可生产的系统能力。
  \end{itemize}
  \item \textbf{\href{https://huggingface.co/papers/2605.00658}{UniVidX：基于扩散先验的通用视频生成统一多模态框架}}（\href{https://arxiv.org/abs/2605.00658}{2605.00658}；每日论文日期：2026-05-04；点赞数：81）
  \begin{itemize}[leftmargin=1.5em]
    \item \textbf{动机}：把看图、懂视频、生成、编辑、搜索和交互纳入统一接口，并提升长程一致性与可控性。
    \item \textbf{背景}：多模态模型正在从单一理解或单一生成任务走向统一系统，但跨模态一致性、可控编辑和长视频稳定性仍是难点。
    \item \textbf{方法}：论文通常通过统一架构、扩散先验、奖励驱动编辑、视频数据规模化或跨模态条件机制，提升生成和理解的一致性。本文题目中的核心对象是“UniVidX：基于扩散先验的通用视频生成统一多模态框架”，可把它理解为该方向的一种具体实现或问题重述。
    \item \textbf{洞察}：生成模型的核心竞争正从单帧质量转向可控性、一致性和工作流可用性。
    \item \textbf{结论}：它的结论价值在于把多模态生成从视觉质量竞争推向可控、可交互、可生产的系统能力。
  \end{itemize}
  \item \textbf{\href{https://huggingface.co/papers/2604.23774}{Prox-E：基于基本图元抽象的细粒度 3D 形状编辑}}（\href{https://arxiv.org/abs/2604.23774}{2604.23774}；每日论文日期：2026-05-04；点赞数：16）
  \begin{itemize}[leftmargin=1.5em]
    \item \textbf{动机}：把看图、懂视频、生成、编辑、搜索和交互纳入统一接口，并提升长程一致性与可控性。
    \item \textbf{背景}：多模态模型正在从单一理解或单一生成任务走向统一系统，但跨模态一致性、可控编辑和长视频稳定性仍是难点。
    \item \textbf{方法}：论文通常通过统一架构、扩散先验、奖励驱动编辑、视频数据规模化或跨模态条件机制，提升生成和理解的一致性。本文题目中的核心对象是“Prox-E：基于基本图元抽象的细粒度 3D 形状编辑”，可把它理解为该方向的一种具体实现或问题重述。
    \item \textbf{洞察}：真正有用的世界模型必须同时回答状态会如何变化、动作是否可行以及失败风险在哪里。
    \item \textbf{结论}：它的结论价值在于把多模态生成从视觉质量竞争推向可控、可交互、可生产的系统能力。
  \end{itemize}
  \item \textbf{\href{https://huggingface.co/papers/2604.28190}{面向视觉生成的表征 Fréchet 损失}}（\href{https://arxiv.org/abs/2604.28190}{2604.28190}；每日论文日期：2026-05-01；点赞数：28）
  \begin{itemize}[leftmargin=1.5em]
    \item \textbf{动机}：把看图、懂视频、生成、编辑、搜索和交互纳入统一接口，并提升长程一致性与可控性。
    \item \textbf{背景}：多模态模型正在从单一理解或单一生成任务走向统一系统，但跨模态一致性、可控编辑和长视频稳定性仍是难点。
    \item \textbf{方法}：论文通常通过统一架构、扩散先验、奖励驱动编辑、视频数据规模化或跨模态条件机制，提升生成和理解的一致性。本文题目中的核心对象是“面向视觉生成的表征 Fréchet 损失”，可把它理解为该方向的一种具体实现或问题重述。
    \item \textbf{洞察}：生成模型的核心竞争正从单帧质量转向可控性、一致性和工作流可用性。
    \item \textbf{结论}：它的结论价值在于把多模态生成从视觉质量竞争推向可控、可交互、可生产的系统能力。
  \end{itemize}
\end{enumerate}

\subsubsection{类别趋势}
趋势：理解和生成边界继续消失，旗舰模型走向任意模态到任意模态，垂直模型则围绕编辑、视频叙事、人类动作、搜索和领域审美构建闭环。

\subsubsection{下一阶段预测}
预测：竞争壁垒会从主干网络参数规模转向数据治理、工作流、可控性、评测闭环和领域化奖励。

\subsection{基础架构、表示学习与推理效率}

本类关注潜空间、词元、注意力、预填充、模型合并和超深网络稳定性等底层问题。

\subsubsection{论文卡片}
\begin{enumerate}[leftmargin=1.5em]
  \item \textbf{\href{https://huggingface.co/papers/2605.06546}{用词元叠加提升预训练效率}}（\href{https://arxiv.org/abs/2605.06546}{2605.06546}；每日论文日期：2026-05-13；点赞数：30）
  \begin{itemize}[leftmargin=1.5em]
    \item \textbf{动机}：突破长上下文、长视频、深层网络和大规模生成的训练/推理成本墙。
    \item \textbf{背景}：模型规模、上下文长度和生成时长继续增长，使注意力、键值缓存、潜表示、流水线训练和模型合并成为基础设施问题。
    \item \textbf{方法}：论文通常通过潜空间建模、残差/注意力结构、键值缓存压缩、预填充优化、流水线并行或模型合并规律，降低成本并提升可扩展性。本文题目中的核心对象是“用词元叠加提升预训练效率”，可把它理解为该方向的一种具体实现或问题重述。
    \item \textbf{洞察}：这篇工作的关键不在单点技巧，而在它把数据、模型、推理、奖励与评测中的某个环节重新组织进系统闭环。
    \item \textbf{结论}：它的结论价值在于为长上下文、长视频和大规模生成提供更现实的计算路径。
  \end{itemize}
  \item \textbf{\href{https://huggingface.co/papers/2605.10780}{超越最后一层：用于视觉词元化的多层表征融合}}（\href{https://arxiv.org/abs/2605.10780}{2605.10780}；每日论文日期：2026-05-13；点赞数：28）
  \begin{itemize}[leftmargin=1.5em]
    \item \textbf{动机}：突破长上下文、长视频、深层网络和大规模生成的训练/推理成本墙。
    \item \textbf{背景}：模型规模、上下文长度和生成时长继续增长，使注意力、键值缓存、潜表示、流水线训练和模型合并成为基础设施问题。
    \item \textbf{方法}：论文通常通过潜空间建模、残差/注意力结构、键值缓存压缩、预填充优化、流水线并行或模型合并规律，降低成本并提升可扩展性。本文题目中的核心对象是“超越最后一层：用于视觉词元化的多层表征融合”，可把它理解为该方向的一种具体实现或问题重述。
    \item \textbf{洞察}：生成模型的核心竞争正从单帧质量转向可控性、一致性和工作流可用性。
    \item \textbf{结论}：它的结论价值在于为长上下文、长视频和大规模生成提供更现实的计算路径。
  \end{itemize}
  \item \textbf{\href{https://huggingface.co/papers/2605.12013}{L2P：释放像素生成中的潜在能力}}（\href{https://arxiv.org/abs/2605.12013}{2605.12013}；每日论文日期：2026-05-13；点赞数：23）
  \begin{itemize}[leftmargin=1.5em]
    \item \textbf{动机}：突破长上下文、长视频、深层网络和大规模生成的训练/推理成本墙。
    \item \textbf{背景}：模型规模、上下文长度和生成时长继续增长，使注意力、键值缓存、潜表示、流水线训练和模型合并成为基础设施问题。
    \item \textbf{方法}：论文通常通过潜空间建模、残差/注意力结构、键值缓存压缩、预填充优化、流水线并行或模型合并规律，降低成本并提升可扩展性。本文题目中的核心对象是“L2P：释放像素生成中的潜在能力”，可把它理解为该方向的一种具体实现或问题重述。
    \item \textbf{洞察}：生成模型的核心竞争正从单帧质量转向可控性、一致性和工作流可用性。
    \item \textbf{结论}：它的结论价值在于为长上下文、长视频和大规模生成提供更现实的计算路径。
  \end{itemize}
  \item \textbf{\href{https://huggingface.co/papers/2605.07721}{记忆高效的循环 Transformer：在循环语言模型中解耦计算与记忆}}（\href{https://arxiv.org/abs/2605.07721}{2605.07721}；每日论文日期：2026-05-12；点赞数：23）
  \begin{itemize}[leftmargin=1.5em]
    \item \textbf{动机}：突破长上下文、长视频、深层网络和大规模生成的训练/推理成本墙。
    \item \textbf{背景}：模型规模、上下文长度和生成时长继续增长，使注意力、键值缓存、潜表示、流水线训练和模型合并成为基础设施问题。
    \item \textbf{方法}：论文通常通过潜空间建模、残差/注意力结构、键值缓存压缩、预填充优化、流水线并行或模型合并规律，降低成本并提升可扩展性。本文题目中的核心对象是“记忆高效的循环 Transformer：在循环语言模型中解耦计算与记忆”，可把它理解为该方向的一种具体实现或问题重述。
    \item \textbf{洞察}：长期记忆不是简单追加上下文，而是需要压缩、权限、隐私和可更新状态共同设计。
    \item \textbf{结论}：它的结论价值在于为长上下文、长视频和大规模生成提供更现实的计算路径。
  \end{itemize}
  \item \textbf{\href{https://huggingface.co/papers/2605.09877}{键值均值机制}}（\href{https://arxiv.org/abs/2605.09877}{2605.09877}；每日论文日期：2026-05-12；点赞数：19）
  \begin{itemize}[leftmargin=1.5em]
    \item \textbf{动机}：突破长上下文、长视频、深层网络和大规模生成的训练/推理成本墙。
    \item \textbf{背景}：模型规模、上下文长度和生成时长继续增长，使注意力、键值缓存、潜表示、流水线训练和模型合并成为基础设施问题。
    \item \textbf{方法}：论文通常通过潜空间建模、残差/注意力结构、键值缓存压缩、预填充优化、流水线并行或模型合并规律，降低成本并提升可扩展性。本文题目中的核心对象是“键值均值机制”，可把它理解为该方向的一种具体实现或问题重述。
    \item \textbf{洞察}：这篇工作的关键不在单点技巧，而在它把数据、模型、推理、奖励与评测中的某个环节重新组织进系统闭环。
    \item \textbf{结论}：它的结论价值在于为长上下文、长视频和大规模生成提供更现实的计算路径。
  \end{itemize}
  \item \textbf{\href{https://huggingface.co/papers/2605.06169}{均值模式尖叫：面向千层扩散 Transformer 的均值—方差分裂残差}}（\href{https://arxiv.org/abs/2605.06169}{2605.06169}；每日论文日期：2026-05-11；点赞数：119）
  \begin{itemize}[leftmargin=1.5em]
    \item \textbf{动机}：突破长上下文、长视频、深层网络和大规模生成的训练/推理成本墙。
    \item \textbf{背景}：模型规模、上下文长度和生成时长继续增长，使注意力、键值缓存、潜表示、流水线训练和模型合并成为基础设施问题。
    \item \textbf{方法}：论文通常通过潜空间建模、残差/注意力结构、键值缓存压缩、预填充优化、流水线并行或模型合并规律，降低成本并提升可扩展性。本文题目中的核心对象是“均值模式尖叫：面向千层扩散 Transformer 的均值—方差分裂残差”，可把它理解为该方向的一种具体实现或问题重述。
    \item \textbf{洞察}：这篇工作的关键不在单点技巧，而在它把数据、模型、推理、奖励与评测中的某个环节重新组织进系统闭环。
    \item \textbf{结论}：它的结论价值在于为长上下文、长视频和大规模生成提供更现实的计算路径。
  \end{itemize}
  \item \textbf{\href{https://huggingface.co/papers/2605.07825}{各向异性的模态对齐}}（\href{https://arxiv.org/abs/2605.07825}{2605.07825}；每日论文日期：2026-05-11；点赞数：26）
  \begin{itemize}[leftmargin=1.5em]
    \item \textbf{动机}：突破长上下文、长视频、深层网络和大规模生成的训练/推理成本墙。
    \item \textbf{背景}：模型规模、上下文长度和生成时长继续增长，使注意力、键值缓存、潜表示、流水线训练和模型合并成为基础设施问题。
    \item \textbf{方法}：论文通常通过潜空间建模、残差/注意力结构、键值缓存压缩、预填充优化、流水线并行或模型合并规律，降低成本并提升可扩展性。本文题目中的核心对象是“各向异性的模态对齐”，可把它理解为该方向的一种具体实现或问题重述。
    \item \textbf{洞察}：这篇工作的关键不在单点技巧，而在它把数据、模型、推理、奖励与评测中的某个环节重新组织进系统闭环。
    \item \textbf{结论}：它的结论价值在于为长上下文、长视频和大规模生成提供更现实的计算路径。
  \end{itemize}
  \item \textbf{\href{https://huggingface.co/papers/2605.07748}{TextLDM：基于连续潜空间扩散的语言建模}}（\href{https://arxiv.org/abs/2605.07748}{2605.07748}；每日论文日期：2026-05-11；点赞数：22）
  \begin{itemize}[leftmargin=1.5em]
    \item \textbf{动机}：突破长上下文、长视频、深层网络和大规模生成的训练/推理成本墙。
    \item \textbf{背景}：模型规模、上下文长度和生成时长继续增长，使注意力、键值缓存、潜表示、流水线训练和模型合并成为基础设施问题。
    \item \textbf{方法}：论文通常通过潜空间建模、残差/注意力结构、键值缓存压缩、预填充优化、流水线并行或模型合并规律，降低成本并提升可扩展性。本文题目中的核心对象是“TextLDM：基于连续潜空间扩散的语言建模”，可把它理解为该方向的一种具体实现或问题重述。
    \item \textbf{洞察}：这篇工作的关键不在单点技巧，而在它把数据、模型、推理、奖励与评测中的某个环节重新组织进系统闭环。
    \item \textbf{结论}：它的结论价值在于为长上下文、长视频和大规模生成提供更现实的计算路径。
  \end{itemize}
  \item \textbf{\href{https://huggingface.co/papers/2605.07755}{通过误差控制动力学重新思考循环模型中的状态跟踪}}（\href{https://arxiv.org/abs/2605.07755}{2605.07755}；每日论文日期：2026-05-11；点赞数：21）
  \begin{itemize}[leftmargin=1.5em]
    \item \textbf{动机}：突破长上下文、长视频、深层网络和大规模生成的训练/推理成本墙。
    \item \textbf{背景}：模型规模、上下文长度和生成时长继续增长，使注意力、键值缓存、潜表示、流水线训练和模型合并成为基础设施问题。
    \item \textbf{方法}：论文通常通过潜空间建模、残差/注意力结构、键值缓存压缩、预填充优化、流水线并行或模型合并规律，降低成本并提升可扩展性。本文题目中的核心对象是“通过误差控制动力学重新思考循环模型中的状态跟踪”，可把它理解为该方向的一种具体实现或问题重述。
    \item \textbf{洞察}：这篇工作的关键不在单点技巧，而在它把数据、模型、推理、奖励与评测中的某个环节重新组织进系统闭环。
    \item \textbf{结论}：它的结论价值在于为长上下文、长视频和大规模生成提供更现实的计算路径。
  \end{itemize}
  \item \textbf{\href{https://huggingface.co/papers/2605.06221}{UniPrefill：通过块级动态稀疏实现通用长上下文预填充加速}}（\href{https://arxiv.org/abs/2605.06221}{2605.06221}；每日论文日期：2026-05-11；点赞数：20）
  \begin{itemize}[leftmargin=1.5em]
    \item \textbf{动机}：突破长上下文、长视频、深层网络和大规模生成的训练/推理成本墙。
    \item \textbf{背景}：模型规模、上下文长度和生成时长继续增长，使注意力、键值缓存、潜表示、流水线训练和模型合并成为基础设施问题。
    \item \textbf{方法}：论文通常通过潜空间建模、残差/注意力结构、键值缓存压缩、预填充优化、流水线并行或模型合并规律，降低成本并提升可扩展性。本文题目中的核心对象是“UniPrefill：通过块级动态稀疏实现通用长上下文预填充加速”，可把它理解为该方向的一种具体实现或问题重述。
    \item \textbf{洞察}：这篇工作的关键不在单点技巧，而在它把数据、模型、推理、奖励与评测中的某个环节重新组织进系统闭环。
    \item \textbf{结论}：它的结论价值在于为长上下文、长视频和大规模生成提供更现实的计算路径。
  \end{itemize}
  \item \textbf{\href{https://huggingface.co/papers/2605.06548}{连续潜空间扩散语言模型}}（\href{https://arxiv.org/abs/2605.06548}{2605.06548}；每日论文日期：2026-05-08；点赞数：70）
  \begin{itemize}[leftmargin=1.5em]
    \item \textbf{动机}：突破长上下文、长视频、深层网络和大规模生成的训练/推理成本墙。
    \item \textbf{背景}：模型规模、上下文长度和生成时长继续增长，使注意力、键值缓存、潜表示、流水线训练和模型合并成为基础设施问题。
    \item \textbf{方法}：论文通常通过潜空间建模、残差/注意力结构、键值缓存压缩、预填充优化、流水线并行或模型合并规律，降低成本并提升可扩展性。本文题目中的核心对象是“连续潜空间扩散语言模型”，可把它理解为该方向的一种具体实现或问题重述。
    \item \textbf{洞察}：这篇工作的关键不在单点技巧，而在它把数据、模型、推理、奖励与评测中的某个环节重新组织进系统闭环。
    \item \textbf{结论}：它的结论价值在于为长上下文、长视频和大规模生成提供更现实的计算路径。
  \end{itemize}
  \item \textbf{\href{https://huggingface.co/papers/2605.04569}{通过上下文稀疏注意力实现快速统一视频编辑}}（\href{https://arxiv.org/abs/2605.04569}{2605.04569}；每日论文日期：2026-05-07；点赞数：17）
  \begin{itemize}[leftmargin=1.5em]
    \item \textbf{动机}：突破长上下文、长视频、深层网络和大规模生成的训练/推理成本墙。
    \item \textbf{背景}：模型规模、上下文长度和生成时长继续增长，使注意力、键值缓存、潜表示、流水线训练和模型合并成为基础设施问题。
    \item \textbf{方法}：论文通常通过潜空间建模、残差/注意力结构、键值缓存压缩、预填充优化、流水线并行或模型合并规律，降低成本并提升可扩展性。本文题目中的核心对象是“通过上下文稀疏注意力实现快速统一视频编辑”，可把它理解为该方向的一种具体实现或问题重述。
    \item \textbf{洞察}：生成模型的核心竞争正从单帧质量转向可控性、一致性和工作流可用性。
    \item \textbf{结论}：它的结论价值在于为长上下文、长视频和大规模生成提供更现实的计算路径。
  \end{itemize}
  \item \textbf{\href{https://huggingface.co/papers/2605.00809}{让视觉 Transformer 说话：生成式语言—图像预训练}}（\href{https://arxiv.org/abs/2605.00809}{2605.00809}；每日论文日期：2026-05-04；点赞数：32）
  \begin{itemize}[leftmargin=1.5em]
    \item \textbf{动机}：突破长上下文、长视频、深层网络和大规模生成的训练/推理成本墙。
    \item \textbf{背景}：模型规模、上下文长度和生成时长继续增长，使注意力、键值缓存、潜表示、流水线训练和模型合并成为基础设施问题。
    \item \textbf{方法}：论文通常通过潜空间建模、残差/注意力结构、键值缓存压缩、预填充优化、流水线并行或模型合并规律，降低成本并提升可扩展性。本文题目中的核心对象是“让视觉 Transformer 说话：生成式语言—图像预训练”，可把它理解为该方向的一种具体实现或问题重述。
    \item \textbf{洞察}：生成模型的核心竞争正从单帧质量转向可控性、一致性和工作流可用性。
    \item \textbf{结论}：它的结论价值在于为长上下文、长视频和大规模生成提供更现实的计算路径。
  \end{itemize}
  \item \textbf{\href{https://huggingface.co/papers/2605.00553}{Stable-GFlowNet：通过对比轨迹平衡实现多样且鲁棒的大模型红队}}（\href{https://arxiv.org/abs/2605.00553}{2605.00553}；每日论文日期：2026-05-04；点赞数：17）
  \begin{itemize}[leftmargin=1.5em]
    \item \textbf{动机}：突破长上下文、长视频、深层网络和大规模生成的训练/推理成本墙。
    \item \textbf{背景}：模型规模、上下文长度和生成时长继续增长，使注意力、键值缓存、潜表示、流水线训练和模型合并成为基础设施问题。
    \item \textbf{方法}：论文通常通过潜空间建模、残差/注意力结构、键值缓存压缩、预填充优化、流水线并行或模型合并规律，降低成本并提升可扩展性。本文题目中的核心对象是“Stable-GFlowNet：通过对比轨迹平衡实现多样且鲁棒的大模型红队”，可把它理解为该方向的一种具体实现或问题重述。
    \item \textbf{洞察}：好的评测应当像诊断仪一样定位失败机制，而不是只给一个总分。
    \item \textbf{结论}：它的结论价值在于为长上下文、长视频和大规模生成提供更现实的计算路径。
  \end{itemize}
  \item \textbf{\href{https://huggingface.co/papers/2604.27085}{RoundPipe：在多张消费级显卡上进行高效训练}}（\href{https://arxiv.org/abs/2604.27085}{2604.27085}；每日论文日期：2026-05-01；点赞数：40）
  \begin{itemize}[leftmargin=1.5em]
    \item \textbf{动机}：突破长上下文、长视频、深层网络和大规模生成的训练/推理成本墙。
    \item \textbf{背景}：模型规模、上下文长度和生成时长继续增长，使注意力、键值缓存、潜表示、流水线训练和模型合并成为基础设施问题。
    \item \textbf{方法}：论文通常通过潜空间建模、残差/注意力结构、键值缓存压缩、预填充优化、流水线并行或模型合并规律，降低成本并提升可扩展性。本文题目中的核心对象是“RoundPipe：在多张消费级显卡上进行高效训练”，可把它理解为该方向的一种具体实现或问题重述。
    \item \textbf{洞察}：这篇工作的关键不在单点技巧，而在它把数据、模型、推理、奖励与评测中的某个环节重新组织进系统闭环。
    \item \textbf{结论}：它的结论价值在于为长上下文、长视频和大规模生成提供更现实的计算路径。
  \end{itemize}
\end{enumerate}

\subsubsection{类别趋势}
趋势：底层研究集中在更深、更长、更连续、更省；连续潜空间、键值缓存压缩、动态稀疏、循环 Transformer 与模型合并都在服务长上下文和长视频。

\subsubsection{下一阶段预测}
预测：系统—模型协同会成为主旋律；连续潜表示与离散词元会长期共存，推理效率论文会越来越直接决定智能体和生成系统能否部署。

\subsection{评测、安全、隐私与可信科学工具}

本类关注基准、隐私、红队、可信推理和科学工具如何为能力增长提供约束。

\subsubsection{论文卡片}
\begin{enumerate}[leftmargin=1.5em]
  \item \textbf{\href{https://huggingface.co/papers/2605.09063}{Soohak：数学家策划的研究级数学能力评测基准}}（\href{https://arxiv.org/abs/2605.09063}{2605.09063}；每日论文日期：2026-05-12；点赞数：71）
  \begin{itemize}[leftmargin=1.5em]
    \item \textbf{动机}：把模型能力放进可审计、可复现、可治理的评测和安全框架中。
    \item \textbf{背景}：能力增长带来幻觉、隐私泄漏、红队攻击、推理不可复现和科研工具可信度问题，需要更贴近真实环境的诊断式评测。
    \item \textbf{方法}：论文通常通过新基准、红队框架、隐私机制、幻觉诊断、数学/世界推理测试或科学文档工具，刻画模型失败边界。本文题目中的核心对象是“Soohak：数学家策划的研究级数学能力评测基准”，可把它理解为该方向的一种具体实现或问题重述。
    \item \textbf{洞察}：搜索能力的关键不只是召回相似文本，而是智能体如何与语料、工具和反馈形成闭环。
    \item \textbf{结论}：它的结论价值在于让模型能力增长能够被更细粒度地诊断、约束和复现。
  \end{itemize}
  \item \textbf{\href{https://huggingface.co/papers/2605.10341}{PaperFit：视觉闭环的科学文档排版优化}}（\href{https://arxiv.org/abs/2605.10341}{2605.10341}；每日论文日期：2026-05-12；点赞数：30）
  \begin{itemize}[leftmargin=1.5em]
    \item \textbf{动机}：把模型能力放进可审计、可复现、可治理的评测和安全框架中。
    \item \textbf{背景}：能力增长带来幻觉、隐私泄漏、红队攻击、推理不可复现和科研工具可信度问题，需要更贴近真实环境的诊断式评测。
    \item \textbf{方法}：论文通常通过新基准、红队框架、隐私机制、幻觉诊断、数学/世界推理测试或科学文档工具，刻画模型失败边界。本文题目中的核心对象是“PaperFit：视觉闭环的科学文档排版优化”，可把它理解为该方向的一种具体实现或问题重述。
    \item \textbf{洞察}：这篇工作的关键不在单点技巧，而在它把数据、模型、推理、奖励与评测中的某个环节重新组织进系统闭环。
    \item \textbf{结论}：它的结论价值在于让模型能力增长能够被更细粒度地诊断、约束和复现。
  \end{itemize}
  \item \textbf{\href{https://huggingface.co/papers/2605.04808}{DTap：面向 AI 智能体的可控交互式红队平台}}（\href{https://arxiv.org/abs/2605.04808}{2605.04808}；每日论文日期：2026-05-11；点赞数：19）
  \begin{itemize}[leftmargin=1.5em]
    \item \textbf{动机}：把模型能力放进可审计、可复现、可治理的评测和安全框架中。
    \item \textbf{背景}：能力增长带来幻觉、隐私泄漏、红队攻击、推理不可复现和科研工具可信度问题，需要更贴近真实环境的诊断式评测。
    \item \textbf{方法}：论文通常通过新基准、红队框架、隐私机制、幻觉诊断、数学/世界推理测试或科学文档工具，刻画模型失败边界。本文题目中的核心对象是“DTap：面向 AI 智能体的可控交互式红队平台”，可把它理解为该方向的一种具体实现或问题重述。
    \item \textbf{洞察}：好的评测应当像诊断仪一样定位失败机制，而不是只给一个总分。
    \item \textbf{结论}：它的结论价值在于让模型能力增长能够被更细粒度地诊断、约束和复现。
  \end{itemize}
  \item \textbf{\href{https://huggingface.co/papers/2605.04523}{RaguTeam：面向忠实多轮回答生成的裁判编排大模型集成}}（\href{https://arxiv.org/abs/2605.04523}{2605.04523}；每日论文日期：2026-05-08；点赞数：42）
  \begin{itemize}[leftmargin=1.5em]
    \item \textbf{动机}：把模型能力放进可审计、可复现、可治理的评测和安全框架中。
    \item \textbf{背景}：能力增长带来幻觉、隐私泄漏、红队攻击、推理不可复现和科研工具可信度问题，需要更贴近真实环境的诊断式评测。
    \item \textbf{方法}：论文通常通过新基准、红队框架、隐私机制、幻觉诊断、数学/世界推理测试或科学文档工具，刻画模型失败边界。本文题目中的核心对象是“RaguTeam：面向忠实多轮回答生成的裁判编排大模型集成”，可把它理解为该方向的一种具体实现或问题重述。
    \item \textbf{洞察}：生成模型的核心竞争正从单帧质量转向可控性、一致性和工作流可用性。
    \item \textbf{结论}：它的结论价值在于让模型能力增长能够被更细粒度地诊断、约束和复现。
  \end{itemize}
  \item \textbf{\href{https://huggingface.co/papers/2605.02910}{CreativityBench：通过基于可供性的工具再利用评估智能体创造性推理}}（\href{https://arxiv.org/abs/2605.02910}{2605.02910}；每日论文日期：2026-05-07；点赞数：21）
  \begin{itemize}[leftmargin=1.5em]
    \item \textbf{动机}：把模型能力放进可审计、可复现、可治理的评测和安全框架中。
    \item \textbf{背景}：能力增长带来幻觉、隐私泄漏、红队攻击、推理不可复现和科研工具可信度问题，需要更贴近真实环境的诊断式评测。
    \item \textbf{方法}：论文通常通过新基准、红队框架、隐私机制、幻觉诊断、数学/世界推理测试或科学文档工具，刻画模型失败边界。本文题目中的核心对象是“CreativityBench：通过基于可供性的工具再利用评估智能体创造性推理”，可把它理解为该方向的一种具体实现或问题重述。
    \item \textbf{洞察}：好的评测应当像诊断仪一样定位失败机制，而不是只给一个总分。
    \item \textbf{结论}：它的结论价值在于让模型能力增长能够被更细粒度地诊断、约束和复现。
  \end{itemize}
  \item \textbf{\href{https://huggingface.co/papers/2605.01428}{幻觉削弱信任：元认知是一条出路}}（\href{https://arxiv.org/abs/2605.01428}{2605.01428}；每日论文日期：2026-05-05；点赞数：22）
  \begin{itemize}[leftmargin=1.5em]
    \item \textbf{动机}：把模型能力放进可审计、可复现、可治理的评测和安全框架中。
    \item \textbf{背景}：能力增长带来幻觉、隐私泄漏、红队攻击、推理不可复现和科研工具可信度问题，需要更贴近真实环境的诊断式评测。
    \item \textbf{方法}：论文通常通过新基准、红队框架、隐私机制、幻觉诊断、数学/世界推理测试或科学文档工具，刻画模型失败边界。本文题目中的核心对象是“幻觉削弱信任：元认知是一条出路”，可把它理解为该方向的一种具体实现或问题重述。
    \item \textbf{洞察}：好的评测应当像诊断仪一样定位失败机制，而不是只给一个总分。
    \item \textbf{结论}：它的结论价值在于让模型能力增长能够被更细粒度地诊断、约束和复现。
  \end{itemize}
\end{enumerate}

\subsubsection{类别趋势}
趋势：评测从静态题库走向动态压力测试、真实工作流、世界状态预测、人类对齐评价细则、可控红队和隐私记忆治理。

\subsubsection{下一阶段预测}
预测：基准会从排行榜变成诊断仪，不仅给分，还定位失败机制、生成反例并提供可训练反馈；来源追踪和隐私会成为智能体记忆的硬约束。

\subsection{优先深读建议}
\begin{enumerate}[leftmargin=1.5em]
  \item \textbf{MolmoAct2 / 世界行动模型}：优先进入具身智能和世界行动模型主线。
  \item \textbf{Stream-R1 / Flow-OPD / MARBLE / AlphaGRPO}：优先进入生成模型后训练与奖励结构工程主线。
  \item \textbf{SenseNova-U1 / Qwen-Image-2.0 / UniVidX}：优先进入统一多模态理解生成主线。
  \item \textbf{OpenSearch-VL / OpenSeeker-v2 / Skill1 / MemPrivacy / δ-mem}：优先进入智能体搜索、技能演化与记忆治理主线。
  \item \textbf{WorldReasonBench / DTap / Soohak / Claw-Eval-Live}：优先进入评测、安全和诊断型基准主线。
\end{enumerate}

\subsection{资料来源}
\begin{itemize}[leftmargin=1.5em]
  \item \textbf{Hugging Face 月度入口}：\href{https://huggingface.co/papers/month/2026-05}{2026 年 5 月每日论文月度页面}。
  \item \textbf{Hugging Face 每日论文首页}：\href{https://huggingface.co/papers}{每日论文首页}。
  \item \textbf{arXiv}：每个论文卡片均链接对应 arXiv 摘要页面。
\end{itemize}
